\documentclass{article}
\usepackage[margin=1in]{geometry}
\usepackage{setspace}
\usepackage{graphicx}
\usepackage{amsmath}
\usepackage{amssymb}
\usepackage{booktabs}
\usepackage{multirow}
\usepackage{tikz}
\usetikzlibrary{arrows.meta}
\usepackage[colorlinks=true,linkcolor=blue,citecolor=blue,urlcolor=blue]{hyperref}  % load last

\numberwithin{equation}{section}   % equations numbered (2.13), where "2" is the section number

\def\DDS{{\rm DDS}}
\def\DD{{\rm DD}}
\def\TDS{{\rm TDS}}
\def\LDS{{\rm LDS}}
\def\LS{{\rm LS}}
\def\LAG{{\rm LAG}}
\def\nn{\nonumber}

% These notes are split across three PDFs, compiled independently from
% notes/dedispersion.tex, notes/detrending.tex, notes/variance_map.tex.
% Everything above this comment is deliberately IDENTICAL in all three sources,
% so that a section can be moved between them without touching the preamble.
%
% A \ref cannot cross documents, so a pointer into a companion PDF names its
% section in prose instead: \ddnotes{Subband search} renders as
% [dedispersion.pdf, "Subband search"]. Each source defines the macros for the
% other two documents (either may be unused at any given moment). The argument
% must reproduce a section or subsection title of the companion verbatim.
\newcommand{\ddnotes}[1]{{\em dedispersion.pdf}, ``#1''}
\newcommand{\dtnotes}[1]{{\em detrending.pdf}, ``#1''}

\title{Notes on the ``pirate'' dedisperser: the variance map (in progress)}
\author{Kendrick Smith}

\begin{document}

\maketitle

\tableofcontents

\setstretch{1.05}   % spacing between lines
\setlength{\parskip}{2.5pt}  % spacing between paragraphs

\clearpage

\section{The ``variance map'' (in progress)}
\label{sec:variance_map}

This section is work in progress, and its relation to the rest of the real-time
pipeline won't be clear yet!

\subsection{Setup}
\label{ssec:variance_map_setup}

Fix a {\tt DedispersionTree} and use notation from \ddnotes{Tree dedispersion}:
\begin{align}
r &= \mbox{tree rank} \\
R &= \mbox{pf\_rank} \\
M &= \mbox{number of multiplets} \\
P &= \mbox{number of peak-finding profiles} \\
\gamma &= \mbox{downsampling level (i.e.\ downsampling factor is $2^\gamma$)}
\end{align}
Consider a linear operator $L$, whose input is an $(N_{\rm freq}, N_{\rm time})$ array
and whose output is a $(2^{r-R}, M, P, N_{\rm time}/2^\gamma)$ array defined by the
following sequence of steps:
\begin{enumerate}
\item Apply a {\tt Detrender2d} to the $(N_{\rm freq}, N_{\rm time})$ array
  in-place with {\bf no mask}. The parameters of the {\tt Detrender2d} are
  part of the problem specification -- we might ``warm up'' with a simple
  case (e.g.\ $n=W=0$).
\item Apply tree dedispersion, including peak-finding convolution for each $0 \le p < P$, 
  but {\bf not} multiplying by a weight array to normalize to ``sigmas'', or coarse-graining.
  Thus, $L$ is the unnormalized dedispersion transform with explicit indices for subbands
  and peak-finding kernel.
\end{enumerate}
Assume that each element of the input $(N_{\rm freq}, N_{\rm time})$ array is an independent
mean-zero Gaussian with known variance $v_F$ which depends on frequency channel
$0 \le F < N_{\rm freq}$ but not time.
Then, after dedispersion reaches steady state, the variance of the output array $y_{dmp}$
depends on $0 \le d < 2^{r-R}$, $0 \le m < M$, $0 \le p < P$, but not time.
We sometimes ``flatten'' the indices $(d,m,p)$ to a single index $0 \le \alpha < 2^{r-R}MP$.
Since $L$ is linear, the output variance $y_\alpha$ is linear in the input variance $v_F$:
\begin{equation}
y_\alpha = \sum_F A_{\alpha F} v_F  \label{eq:variance_map_def}
\end{equation}
where this equation defines the {\bf variance map} $A$, a $(2^{r-R}MP)$-by-$N_{\rm freq}$ matrix.

In its dense matrix representation, the variance map is too expensive to compute (or even store)
at CHIME/CHORD scale.
This section is concerned with the following central, currently unsolved, problem:
can we find a {\bf low-rank} approximation to the variance map $A \approx QW^T$, which can be computed
in a reasonable amount of time?
That is, we are trying to achieve three goals:
\begin{enumerate}
\item We want to minimize the rank of the approximation $A^{\rm approx} \equiv QW^T$.
\item We want the approximation $A^{\rm approx} \approx A^{\rm true}$ to be accurate,
 in the sense that a certain distance function $D(A^{\rm true}, A^{\rm approx})$ is
 small. The choice of distance function is unusual, and is discussed in the next subsection.
\item We want $A^{\rm approx}$ to be computable in a reasonable amount of time.
 In particular, we'd like to avoid computing every matrix element of $A^{\rm true}$,
 since this is cost-prohibitive at CHIME/CHORD scale.
 This aspect of the problem is particularly challenging, so I suggest ``warming up'' by
 exploring subscale examples, where computing $A^{\rm true}$ is feasible, and finding
 a scalable solution later.
\end{enumerate}

\subsection{Distance function}
\label{ssec:distance_function}

Given matrices $A^{\rm true}_{\alpha F}, A^{\rm approx}_{\alpha F}$,
our distance function $D(A^{\rm true}, A^{\rm approx})$ is defined as follows.
First, we define vectors
$y^{\rm true}_{\alpha}$, $y^{\rm approx}_{\alpha}$ by:
\begin{equation}
y^{\rm true}_{\alpha} \equiv \sum_F A^{\rm true}_{\alpha F}
  \hspace{1.5cm}
y^{\rm approx}_{\alpha} \equiv \sum_F A^{\rm approx}_{\alpha F}
\end{equation}
Then we define:
\begin{equation}
D(A^{\rm true}, A^{\rm approx})
 \equiv \begin{cases}
   \infty & \mbox{if $A^{\rm approx}_{\alpha F} < A^{\rm true}_{\alpha F}$ for any $(\alpha,F)$} \\
   N_\alpha^{-1} \sum_\alpha f(y_\alpha^{\rm approx}/y_\alpha^{\rm true})
    & \mbox{otherwise}
\end{cases} 
\label{eq:distance_function}
\end{equation}
where $N_\alpha = \sum_\alpha 1$ is the number of $\alpha$-values, and $f(x)$ is defined by:
\begin{equation}
f(x) \equiv \frac{x-1}{1+(x/10)}
\end{equation}
This is implemented by \verb|class VarMapDistance| (in \verb|pirate_frb/slow_avar/|),
which should be used wherever a distance is reported, so that results are comparable.
Besides $D$ itself, it exposes the quantity $D_0$ (the second branch above, evaluated
unconditionally) and the worst per-element ratio
$\max_{\alpha F} (A^{\rm true}_{\alpha F} / A^{\rm approx}_{\alpha F})$ together with the
$(\alpha,F)$ where it occurs. These are more informative than $D$ alone: if the max ratio
is only a little larger than $1$, then inflating $A^{\rm approx}$ slightly gives a finite
distance close to $D_0$, so such an approximation is nearly usable -- whereas
$D = \infty$ by itself does not distinguish it from a hopeless one.

This unusual distance function has the following properties:
\begin{enumerate}
\item Underestimating any matrix element of $A$ is catastrophic -- the approximation
 is considered infinitely bad.
\item Overestimating is okay, but our goal is to minimize the average level of
 overestimation.
\item Extreme overestimation is fine, as long as it's confined to a small number
 of $\alpha$-values. (If we send $y_\alpha \rightarrow \infty$ for a few $\alpha$
 values, then it has a bounded effect on the distance function.)
\end{enumerate}
The motivation is as follows.
In the real-time search, we normalize dedispersion outputs to ``sigmas'' by multiplying
by $(Av)^{-1/2}$, where $v_F$ is a real-time estimate for the RFI-masked variance.
We search a gigantic number of trials ($\sim\!10^{15}$) for rare events, and we use
a high threshold (9--10$\sigma$) due to the high trial factor.

Underestimating variance is catastrophic, even in corner cases, since it can generate
a large number of false positives in the far tail of the SNR distribution.
The RFI mask is unpredictable and can produce corner cases, such as a small number
of frequency channels contributing.
Therefore, we impose a strict constraint (per matrix element!) that we never underestimate
variance.

In contrast, overestimating variance is undesirable, since it reduces the total event rate,
but this is a statistical concern, not a hard constraint.
Because it's a statistical concern, our cost function focuses on the typical case where
all channels have similar variance, not per-matrix-element analysis.

\subsection{Current status}

\begin{itemize}

\item If we simplify the problem by removing the {\tt Detrender2d}, then
 there is an exact low-rank representation for $A$, and an approximation which
 is fast to compute (see the appendices). This is all implemented and tested.
 I don't know if it's a useful starting point -- it's not clear whether the
 approach can be generalized to include the {\tt Detrender2d}.

\item We have code to compute the $A$-matrix in dense form (\S\ref{ssec:variance_map_bruteforce}),
 given a {\tt DedispersionConfig} and a {\tt Detrender2dParams}. (See \verb|variance_map| in \verb|__main__.py|.)
 This may be a useful ``crutch'' for initial exploration, although in the long run it won't
 scale to CHIME/CHORD. For example, we could compute the full SVD decomposition of the $A$-matrix,
 as a starting point for finding low-rank approximations. It would be a significant milestone
 to develop an algorithm which works well starting from a dense representation of $A$.

\end{itemize}

\subsection{How well can we do starting from a precomputed $A$-matrix?}
\label{ssec:first_milestone}

In this subsection, we'll explore algorithms for finding low-rank approximations,
under the assumption that the dense $A$-matrix has been precomputed and saved to disk.

This subsection is organized as alternating subsubsections which are human-authored
or agent-authored.
The human-authored sections usually propose ideas for the agent to explore, and the
agent-authored sections usually report results (after running a ``campaign'' via the
\verb|ch-research| skill).

\subsubsection{Campaign 1 setup [human-authored]}

Let's start by considering the following question.
Suppose that we've precomputed the variance map $A$, given a {\tt DedispersionConfig} 
and a {\tt Detrender2dParams}, and written it to disk.
What are some interesting algorithms for finding a low-rank approximation with
small $D(A^{\rm true}, A^{\rm approx})$?

An ideal milestone would be to have some code which reads a precomputed variance map $A$,
and returns a few (matrix rank, distance) pairs, to get a sense for how the level of approximation
might depend on the rank of the approximation.

Some practical notes, to get started:
\begin{itemize}

\item To compute an $A$-matrix, use \verb|pirate_frb variance_map <config> <detrender> -o A.asdf|
 (\S\ref{ssec:variance_map_bruteforce}), and to read it back, use
 \verb|slow_avar.read_variance_map()|.

\item To evaluate the distance function, use \verb|class VarMapDistance| (in
 \verb|pirate_frb/slow_avar/|, see \S\ref{ssec:distance_function}). Please use it, rather than
 reimplementing $D$, so that distances reported by different experiments are comparable.
 Besides \verb|D|, report \verb|D0| and \verb|max_r|, since these distinguish an approximation
 which is nearly usable from one which is hopeless.

\item There is no subscale config in \verb|configs/dedispersion/| yet, so part of the task is to
 create a variety of them, exploring the {\tt DedispersionConfig} and {\tt Detrender2dParams}
 parameter space. We are more interested in {\em typical} cases than in adversarial or edge
 cases: the goal is an algorithm that works well on the configs we would actually run, not one
 that survives a worst case we will never encounter.

\end{itemize}

Here are some initial thoughts and directions that may be interesting, but
I encourage you to propose new ideas and directions:

\begin{itemize}

\item I suspect that heuristic algorithms and empirical approaches will
 work best, and that making rigorous guarantees will be hard. (But I'd love to
 be wrong about this!)

\item Let's disable downsampled trees and early triggers. I suspect that
 if we can solve the problem for a single {\tt DedispersionTree}, we'll
 be able to solve the multi-tree case.

\item To start out, let's also disable frequency subbands, so that the
 main parameters to explore are the frequency zones, and the {\tt Detrender2d}.
 Let's also use subscale instances (say rank-10 and $\sim$400 frequency channels).
 Later, we can scale up, and add frequency subbands.

\item Our distance function (\ref{eq:distance_function}) is discontinuous,
 which makes it hard to use optimization-based approaches. It may be useful to
 ``relax'' to a continuous distance function in an intermediate stage. Or, if
 continuity turns out to be very helpful, then we could think about changing
 the definition of the cost function.

\item Suppose we have an accurate approximation $A_1 = Q W^T$ with
 a large rank (say 1000). Can we find an algorithm to search for a lower-rank
 (say 100) representation $A_2$ with small $D(A_1,A_2)$? If such an algorithm
 exists, it could be useful for ``bootstrapping'' low-rank approximations from
 medium-rank ones.

\item Is coarse-graining DM a useful idea? Recall that $0 \le \alpha < (2^{r-R}MP)$
is a ``flattened'' index. Let $0\le \beta < (2^{r-L}NP)$ be a coarser ``flattened''
index obtained by coarse-graining DM. Here, $L \ge R$ is a coarse-graining factor, 
and $N \le M$ is the number of frequency subbands (a multiplet is a subband plus
a fine-grained DM). Define the coarse-grained variance map $\bar A$ by:
\begin{equation}
\bar A_{\beta F} = \max_{\alpha\in\beta} A_{\alpha F}
\label{eq:coarse_graining}
\end{equation}
One approach to controlling rank/cost is to find a low-rank approximation to $\bar A$, 
and ``lift'' it to a low-rank approximation to $A$. (Note that our current ``production''
GPU kernels use coarse-grained weights, which implicitly assumes that this approach is
viable at some level.)

\item If we're having trouble at low DM specifically, then we could inflate variance
estimates at low DM (with little or no change in the rank). Given our unusual distance
function $D$, it may not change the value of $D$ dramatically.

\end{itemize}

\subsubsection{Campaign 1 results [agent-authored]}
\label{sssec:varmap_proposed}

This subsubsection describes an algorithm which came out of an initial exploration of the
milestone above, and the numbers it produces. The short version: the approximation which works is
a {\em one-sided nonnegative factorization at coarse-grained DM}, built by alternating two
linear programs. It is more accurate than any clustering we tried, it turns out to be cheaper to
apply than the alternatives, it never forms the dense $A$, and its output is the array the
production peak-finding kernel already stores.

Throughout, $Q$ and $W$ are the factors of $A^{\rm approx} = QW^T$ as in
\S\ref{ssec:variance_map_setup}: $W$ is $N_{\rm freq}$-by-$K$ and $Q$ is $(2^{r-R}MP)$-by-$K$.
Note the symbol clash with the {\tt Detrender2d} time half-width $W$ of \dtnotes{2-d detrending};
below, the detrender's half-width is always named in words, never as $W$.

\paragraph*{Why not least squares.}
It is worth being explicit about why the obvious tool fails, because it shapes everything else.
The constraint in Eq.\ (\ref{eq:distance_function}) is {\em per matrix element}: a single
$A^{\rm approx}_{\alpha F} < A^{\rm true}_{\alpha F}$ anywhere sends $D$ to infinity. So this is
a {\em one-sided} (upper-bound) approximation problem, not a least-squares problem, and a
truncated SVD is not merely suboptimal but usually inadmissible. A truncation has negative
entries; clipping them to zero still leaves it underestimating $A$ somewhere, and no rescaling
can repair that. On a rank-8 map with a detrender, truncations at rank 8, 16, 32 and 64 are all
inadmissible ($D = \infty$); the first admissible one is rank 128, and it scores $D = 0.10$
where a rank-128 envelope scores $0.054$. Least-squares accuracy is close to irrelevant here,
and at low rank it is actively misleading.

The exception is instructive. With {\em no} {\tt Detrender2d}, $A$ really is close to low rank
-- numerically ${\rm rank}(A) = 54$ to $138$ on our maps -- and a clipped, inflated rank-16 SVD
reaches $D \approx 10^{-6}$, four orders of magnitude better than any clustering at the same
rank. Frequency subbands alone do not break this. It is the {\tt Detrender2d} that destroys the
algebraic structure (\dtnotes{2-d detrending}: it smears every one-hot across its whole spline
zone, which also destroys the exact zeros of $A$), and once it is switched on the SVD is dead at
every rank below $N_{\rm freq}$. The algorithm below is aimed at that case, since it is the one
we run.

\paragraph*{The shape of the approximation: a coarse-grained $Q$.}
Partition the $N_\alpha \equiv 2^{r-R}MP$ output indices $\alpha$ into $N_\beta$ groups, indexed by
$\beta$, and write $\beta(\alpha)$ for the group containing $\alpha$. Give each {\em group} a
single coefficient vector $q_\beta \in \mathbb{R}^K_{\ge 0}$, and let every $\alpha$ in that
group share it:
\begin{equation}
A^{\rm approx}_{\alpha F} = (W q_{\beta(\alpha)})_F
\label{eq:varmap_coarse_assign}
\end{equation}
To fix the shapes: $W$ is $N_{\rm freq}$-by-$K$, so its $K$ columns are frequency vectors;
$q_\beta$ is a column vector of $K$ nonnegative numbers; $W q_\beta$ is
therefore an $N_{\rm freq}$-vector, and $(\cdot)_F$ takes its $F$-th component. In words,
Eq.\ (\ref{eq:varmap_coarse_assign}) says {\em every output row in a group is the same
nonnegative combination of the same $K$ columns of $W$}.

This is the notes' own coarse-graining of \S\ref{ssec:first_milestone} with one generalization:
instead of giving each group a single column of $W$, we let it use a nonnegative
{\em combination} of them.

{\bf Relation to $Q$, and the degrees of freedom.} Collect the $q_\beta$ as the rows of a
$N_\beta$-by-$K$ matrix $\bar Q$. Then the $Q$ of $A^{\rm approx} = QW^T$ is $\bar Q$ with its rows
duplicated across each group,
\begin{equation}
Q = S \bar Q,
\hspace{1.5cm}
S_{\alpha\beta} = \begin{cases} 1 & \mbox{if $\beta(\alpha) = \beta$} \\ 0 & \mbox{otherwise}\end{cases}
\label{eq:varmap_qbar}
\end{equation}
where $S$ is the $N_\alpha$-by-$N_\beta$ one-hot group-membership matrix. So $Q$ formally has
$N_\alpha$ rows, but only $N_\beta$ {\em distinct} ones, and the free parameters are $W$ and
$\bar Q$ alone: $K(N_{\rm freq} + N_\beta)$ of them, against $K(N_{\rm freq} + N_\alpha)$ for an
unrestricted per-row factorization. The rank is $K$ either way.

That row duplication is not a technicality -- it is where the practical win comes from, since
$N_\alpha / N_\beta = 2^{L-R} (M/N)$ is large (233 at CHORD tree 0). At the CHIME tree-0 geometry
below, with $N_\alpha = 524288$, $N_\beta = 8192$, $N_{\rm freq} = 16384$ and $K = 32$: $W$ carries
$5.2\times10^5$ parameters and $\bar Q$ carries $2.6\times10^5$ dense -- or $1.6\times10^4$
nonzeros once $q_\beta$ is capped at two nonzero entries -- against $1.7\times10^7$ for a
per-row $Q$ of the same rank.

The groups are exactly the $\beta$ of \S\ref{ssec:first_milestone}, so that
$N_\beta = 2^{r-L}NP$: coarse-grain DM by a factor
$2^{L-R}$, and never merge across a frequency subband or a peak-finding profile. An ablation
over the three axes says this is not a stylistic choice -- merging subbands costs a factor
$489$ to $3174$ at equal rank, merging profiles a factor $11$ to $175$, whereas merging
{\em fine-grained} DM within a multiplet is nearly free (a factor $1.0$ to $3.9$). The reason
subbands cannot be merged is that they have disjoint frequency support, and this survives the
detrender: the support stops being disjoint but stays concentrated, which is enough.

Define the coarse-grained variance map, as in \S\ref{ssec:first_milestone},
\begin{equation}
\bar A_{\beta F} = \max_{\alpha \in \beta} A_{\alpha F}
\label{eq:varmap_abar}
\end{equation}
Then any $q_\beta$ satisfying $W q_\beta \ge \bar A_{\beta,:}$ elementwise is automatically
admissible for every fine row in the group, since $(Wq_\beta)_F \ge \bar A_{\beta F} \ge
A_{\alpha F}$. {\bf This is the pivot of the whole algorithm}: admissibility becomes a property
of $\bar A$ alone, so the optimization never needs to see a dense row of $A$.

\paragraph*{What the algorithm needs, and why it is cheap.}
Only three things: $\bar A$, the fine row sums $y_\alpha = \sum_F A_{\alpha F}$, and the group
labels (pure geometry). The row sums are needed only to {\em score} the result, not to build it.
Both $\bar A$ and $y$ are computable in a single streaming pass, because the brute-force sweep
of \S\ref{ssec:variance_map_bruteforce} already produces $A$ one input channel (one column) at a
time: max-reduce each column into its groups as it appears, accumulate $y$, and discard the
column. For a CHIME-shaped tree this is $1.0$ GiB of $\bar A$ instead of $64$ GiB of $A$; for
CHORD, $5.4$ GiB instead of $1251$ GiB. Measured against a dense reduction on the $64$ GiB case,
the streamed $\bar A$ is bit-identical, uses $1.9$ GiB of peak RSS against $128.5$ GiB, and is
20\% faster.

\paragraph*{Step 1: constructing $W$ (this is also the initialization).}
We need an initial $W$, i.e.\ $K$ frequency vectors from which every group's row can be covered.
Take the $N_\beta$ rows of $\bar A$, normalize each to unit sum,
\begin{equation}
\bar u_{\beta F} = \bar A_{\beta F} \Big/ \sum_{F'} \bar A_{\beta F'}
\end{equation}
and cluster them by greedy agglomerative merging: start with one cluster per group, and
repeatedly merge the pair whose merge increases the objective least. The column of $W$ for
cluster $c$ is the elementwise max over its members, $W_{Fc} = \max_{\beta \in c} \bar
u_{\beta F}$.

Two things about this are worth understanding rather than just implementing.

{\em Why shapes and not rows.} If a group is given a single column of $W$, it is always
free to rescale it -- $q_\beta = s e_c$ costs nothing in rank -- and the optimal $s$ is
$\max_F \bar A_{\beta F}/W_{Fc}$. Substituting it, the resulting $y^{\rm approx}_\beta /
y^{\rm true}_\beta$ is invariant under $W \to cW$ and does not involve the row's overall
normalization at all. So {\em once rescaling is allowed, only the shapes matter}, and clustering
the raw rows optimizes the wrong thing -- visibly so, since it produces a non-monotone frontier
in $K$. Clustering the unit-sum shapes instead is worth 3--20\% (typically 11--14\%) at no extra
cost, and the argument is stronger for the full algorithm below, whose $q_\beta$ is exactly
scale-invariant in each atom. Note this construction of $W$ needs only $\bar A$ -- not even $y$.

{\em Why this is a good enough starting point.} The initial $W$ is a max envelope, so the
one-hot $q_\beta = e_{c(\beta)}$ is feasible: the alternation starts from an
admissible point with a known finite $D$, and every step below preserves feasibility. The
$W$ also does not have to be optimal, only reasonable: building it from $\bar A$ at the
same coarse-graining used for $Q$ is within 3.6\% of building it from a finer $\bar A$,
and within 3.5\% of building it from the dense $A$. So there is no reason to keep a finer
$\bar A$ around -- the memory footprint is set by $Q$ alone.

\paragraph*{Step 2: the Q-step (optimize $Q$, holding $W$ fixed).}
Let $s_c = \sum_F W_{Fc}$, so that a group's approximate row sum is $s \cdot q_\beta$. Since $f$
in Eq.\ (\ref{eq:distance_function}) is strictly increasing, minimizing $D$ over $q_\beta$ is
exactly minimizing the row sum, and the groups are independent. So the Q-step is, for each
$\beta$ separately, the covering linear program
\begin{equation}
\min_{q_\beta \ge 0} \ s \cdot q_\beta
\hspace{1.2cm}\mbox{subject to}\hspace{1.2cm}
W q_\beta \ \ge \ \bar A_{\beta,:}
\label{eq:varmap_qstep}
\end{equation}
with $K$ variables and $N_{\rm freq}$ constraints. Three remarks. It is {\em exact}, not a
heuristic -- given $W$, no better $Q$ exists. It uses only that $f$ is increasing, so it
is insensitive to the exact shape of the cost function. And it is embarrassingly parallel over
groups. At $N_{\rm freq} = 16384$ it should be solved with constraint generation (start from a
subset of channels, add violated ones, repeat); this is $8.7\times$ faster than handing the
solver all the constraints and agrees with it to $10^{-9}$, and without it the method looks
unaffordable when it is not.

\paragraph*{Step 3: the W-step (optimize $W$, holding $Q$ fixed).}
This one needs a small trick, because the objective depends on $W$ through the column sums while
the constraint depends on $W$ elementwise, so it does not decouple as written. Observe that $f$
is concave:
\begin{equation}
f'(x) = \frac{1.1}{(1+x/10)^2} > 0,
\hspace{1.5cm}
f''(x) < 0
\end{equation}
For a concave $f$ the tangent at the current iterate is a global {\em upper} bound, $f(x) \le
f(x_0) + f'(x_0)(x-x_0)$, so minimizing the tangent is a majorize-minimize step and is guaranteed
not to increase the true objective. The majorizer is linear in $y$, hence linear in $W$: with
weights $w_\alpha = f'(x_\alpha)/y^{\rm true}_\alpha$ evaluated at the current iterate, it is
$\sum_{F,c} W_{Fc}\,(\sum_\alpha w_\alpha Q_{\alpha c})$, which {\em decouples over frequency
channels}. So the W-step is, for each channel $F$ separately,
\begin{equation}
\min_{W_{F,:} \ge 0} \ g \cdot W_{F,:}
\hspace{1.2cm}\mbox{subject to}\hspace{1.2cm}
Q\,W_{F,:} \ \ge \ \bar A_{:,F},
\hspace{1.2cm} g_c \equiv \sum_\alpha w_\alpha Q_{\alpha c}
\label{eq:varmap_wstep}
\end{equation}
with $K$ variables and $N_\beta$ constraints. Again only $\bar A$ is touched.

\paragraph*{Iteration schedule.}
Feasibility is an invariant rather than a hope: after a W-step the current $q_\beta$ is still
feasible for the next Q-step, and after a Q-step the current $W_{F,:}$ is still feasible for the
next W-step. Starting from the admissible envelope of Step 1, every LP in every run we did was
feasible.

In practice, {\bf run three Q-steps and two W-steps and stop}. The first Q-step is 68--94\% of
the total gain; the W-step is worth only a further 7--18\%, and {\em nothing at all} on a map
with no detrender. Monotonicity was checked over 435 consecutive half-steps: the only increases
were $\sim\!10^{-5}$ relative, all attributable to the admissibility repair below.

\paragraph*{Numerical care (this is not optional).}
An LP solution that is ``feasible'' to the solver can still violate $W q \ge \bar A$ enough to
make $D$ infinite. The tolerance to budget for is {\em relative}, but a solver's tolerance is
{\em absolute} -- HiGHS's default primal feasibility tolerance is $\sim\!10^{-7}$, which on a
matrix element of $10^{-12}$ is meaningless. On one subbanded map the W-step set $W_{Fc} = 0$
where $A^{\rm true} > 0$, which no rescaling can repair. Two fixes, both cheap:
\begin{enumerate}
\item {\bf Equilibrate every constraint to a right-hand side of 1}, i.e.\ rewrite $Mx \ge b$ as
 $(M/b)x \ge 1$, dropping the $b=0$ rows (automatically satisfied since $M, x \ge 0$). The
 solver's absolute tolerance is then a uniform relative one. This brought the worst violation
 down to $1.3\times10^{-7}$ -- exactly the solver tolerance -- on every map.
\item {\bf Repair exactly and locally afterwards}, rather than applying a blanket multiplicative
 margin: scale row $\beta$ of $Q$ by $\max(1, \max_F \bar A_{\beta F}/(Wq_\beta)_F)$, and
 similarly per channel after a W-step. Scaling one row of $Q$ touches only that row of the
 product, so the repair cannot break anything else.
\end{enumerate}
With both in place, all 1944 approximations evaluated in the study were admissible with
${\rm max\_r} = 1$ exactly.

\paragraph*{Choosing the coarse-graining factor $L$.}
There is a usable design rule. The {\em floor} $D_{\rm floor}(L)$ -- the plain max envelope of
Eq.\ (\ref{eq:varmap_abar}) lifted back to $A$, with no LP at all -- is a hard lower bound for
the whole coarse-grained-$Q$ family at any rank, and it is cheap to evaluate. Empirically the
achieved distance adds it in quadrature to the distance the same algorithm would reach with
unlimited $L$:
\begin{equation}
D(L) \ \approx \ \sqrt{D_\infty^2 + D_{\rm floor}(L)^2}
\label{eq:varmap_quadrature}
\end{equation}
which held to $-6$\%/$+11$\% over 25 predictions spanning five octaves of $2^{r-L}$. So: pick $L$
such that $D_{\rm floor}(L) \lesssim D_{\rm target}/3$, which requires only a streamed $\bar A$
and one envelope evaluation, and then the LP machinery is worth building.

A second empirical law is useful for extrapolating between configurations. At fixed
$n_{\rm dm,wt} \equiv 2^{r-L}$ (the {\em number} of coarse DM bins), $D$ varies by only
$1.15$--$1.37\times$ over tree ranks 12 to 16, whereas at fixed $L$ it varies by 12--$142\times$.
So $D$ is set by $n_{\rm dm,wt}$, not by $L$, and in the fine regime
\begin{equation}
D \ \approx \ C(r) \, n_{\rm dm,wt}^{-0.85},
\hspace{1.2cm} C \approx 1.3 \ \mbox{to} \ 2.6 \ \mbox{for} \ r = 8 \ \mbox{to} \ 16
\end{equation}
The exponent is $0.85$ and not $1$; assuming $1$ under-predicts $D$ by roughly $2\times$ per
four octaves of extrapolation. Because a fixed \verb|wt_dm_downsampling| is a {\em finer}
fractional DM cut at higher rank, production is better off than any subscale measurement
suggests.

\paragraph*{Results.}
All distances below are from \verb|class VarMapDistance| (\S\ref{ssec:distance_function}) and all
approximations are admissible (${\rm max\_r} = 1$). ``Envelope'' is the best plain clustering we
found, which is the same greedy agglomerative merge in shape space but with one-hot $q_\beta$.

One quantity in the table needs defining first. Write ${\rm nnz}$ for the number of {\em nonzero}
entries of $q_\beta$: how many of the $K$ columns of $W$ a group actually uses. It is a
separate budget from the rank -- $K$ is the number of {\em shared} columns of $W$, ${\rm nnz}$ is
how much of it each group draws on -- and the two enter the cost differently (see below).
The plain envelope is ${\rm nnz} = 1$ by construction, since each group uses one column of $W$.
The Q-step of Eq.\ (\ref{eq:varmap_qstep}) is {\em not} asked to be sparse and returns
${\rm nnz} = 2$ to $3$ anyway. The ``${\rm nnz} \le 2$'' row is the variant where it is
additionally constrained: keep the best admissible subset of size 2 of the LP's own support,
then re-run the alternation with that support frozen.

The cleanest measurement is at the real CHIME tree-0 geometry ($r{=}15$, $N{=}1$, $P{=}16$,
$N_{\rm freq} = 16384$, $N_\alpha = 524288$, dense $|A| = 64$ GiB), coarse-grained at
$L$ corresponding to the shipped \verb|wt_dm_downsampling| $=64$, i.e.\ $N_\beta = 8192$ groups.
The {\tt Detrender2d} is the one used throughout this study: a single spline zone with
$n_\phi = 2$, $n = 2$, time half-width $2$, $T = 2048$, $\eta = 10^{-3}$,
$\epsilon = 3\times10^{-5}$, and a uniform knot vector giving 8 knot intervals across the band
(hence $N_\phi = 10$ spline basis functions), held fixed as $N_{\rm freq}$ varies.
{\bf Note the time half-width is 2, which is the pessimistic end of the range}: it is the
setting that makes the low-DM rows hard, and the distances below would improve by roughly
$2.8\times$ at half-width 8 (see the last bullet of ``what is not yet established'' below).
The table is therefore a conservative operating point, not a tuned one.

\begin{center}
\begin{tabular}{lcccc}
rank $K$ & 8 & 16 & 32 & 64 \\
\hline
(a) envelope (one-hot, ${\rm nnz} = 1$) & 0.0750 & 0.0504 & 0.0357 & 0.0261 \\
(b) {\bf this algorithm} (${\rm nnz}$ unconstrained) & {\bf 0.0412} & {\bf 0.0277} & {\bf 0.0191} & {\bf 0.0136} \\
(c) this algorithm, ${\rm nnz} \le 2$ & 0.0470 & 0.0331 & 0.0235 & 0.0182 \\
\hline
gain at equal rank, (a)/(b) & 1.82$\times$ & 1.82$\times$ & 1.87$\times$ & 1.92$\times$ \\
gain at equal rank, (a)/(c) & 1.60$\times$ & 1.52$\times$ & 1.52$\times$ & 1.43$\times$ \\
\end{tabular}
\end{center}

Every entry is a distance $D$, and all three rows are the same rank $K$: they differ only in how
many columns of $W$ each group may use. Constraining ${\rm nnz}$ to 2 costs 14--34\% in $D$
and is worth doing only when it buys more than that in cost, which at this geometry it does not
-- see below.

For reference, the coarse-graining alone -- one weight per group, i.e.\ rank $N_\beta = 8192$ and
no factorization at all, which is what production currently computes -- gives $D = 0.0043$ on
the same map. That is the floor $D_{\rm floor}(L)$ of Eq.\ (\ref{eq:varmap_quadrature}).

{\bf The trade this exposes is the useful output.} Rank 8192 at $D = 0.0043$ costs
$1.34\times10^8$ multiply-adds to apply; rank 32 with ${\rm nnz}=2$ at $D = 0.0235$ costs
$5.4\times10^5$, i.e.\ {\em 248$\times$ cheaper to apply for 5.4$\times$ more $D$}; rank 8 is
$910\times$ cheaper for $11\times$ more $D$. Which point on that curve is wanted is a science
question -- how much conservative overestimation of the output variance the peak-finder's
threshold can absorb -- not an algorithmic one, and this is the curve on which to answer it.

At equal {\em apply cost} rather than equal rank, the factorization is $2.4$--$3.1\times$ better
in $D$ than the envelope at this geometry. On subscale detrended maps the equal-rank gain is
$1.23$--$1.55\times$, so nothing about the gain was an artefact of small maps -- it grows with
scale. Restricting the algorithm to see only $\bar A$, rather than giving it dense per-row
access, costs $0.93$--$1.06\times$: nothing, and occasionally better, since the coarse constraint
acts as a mild regularizer.

Two further points, which were not obvious to us in advance:

\begin{itemize}

\item {\bf It is cheaper to apply, not just more accurate.} Evaluating $y = Q(W^Tv)$ costs
 $K N_{\rm freq} + N_\beta\,{\rm nnz}$ multiply-adds, where ${\rm nnz}$ is the number of $W$-columns
 vectors a group actually uses (the LP returns 2--3, without being asked). At the CHIME tree-0
 geometry this buys $2.4$--$3.1\times$ in $D$ at equal apply cost. Note the
 saving comes from coarse-graining $Q$: with a {\em per-row} $q_\alpha$ the same factorization is
 {\em worse} than the envelope at equal cost, because the $N_\alpha\,{\rm nnz}$ term dominates.
 Whether to additionally cap ${\rm nnz}$ at 2 depends on which term dominates: cap when
 $N_\beta \gtrsim K N_{\rm freq}/4$, which is true at $N_{\rm freq} = 400$ and false at
 $N_{\rm freq} = 16384$.

\item {\bf The output is the array the kernel already stores.} \verb|PeakFindingKernelParams|'
 weight array has shape $(n_{\rm dm,wt}, n_{t,\rm wt}, P, N)$, i.e.\ one scalar per
 coarse-graining group $\beta$ -- which is precisely Eq.\ (\ref{eq:varmap_coarse_assign}) with
 ${\rm nnz} = 1$. So this algorithm's $Q$ is the object production already has, generalized from
 one entry per group to two or three; ${\rm nnz} = 2$ means the peak-finder applies the existing
 weight multiply twice and adds. A per-{\em row} $Q$ is not expressible in the current
 kernel at all. The reduction is therefore not only cheaper, it is the only one of the two that
 fits the hardware.

\end{itemize}

\paragraph*{Two ideas that did not work.}
Recorded because both look attractive and both cost real effort to rule out.

The {\bf ``give up on the hard rows''} strategy suggested by the saturation of $f$ does not pay.
Abandoning a fraction $q$ of rows (sending them to $y \to \infty$) adds $10q$ to $D$ and removes
the share of $D$ those rows carried, and empirically that share is only $\approx 5q$ -- the
cost is broad, not concentrated. Break-even therefore needs $D \gtrsim 2$, and measured $D$ is
below that everywhere; 71 of 72 direct tests lost. The related suggestion of inflating variance
at low DM should be dropped for the same reason.

{\bf Non-uniform DM cuts} (log, sqrt, inverse-log spacings) are no better than uniform ones.

\paragraph*{What is not yet established.}
\begin{itemize}
\item The high-rank ladder ($r > 12$) is all $N=1$, because no compiled
 {\tt GpuSbDedispersionKernel} exists for $N>1$ at {\tt dd\_rank} 7--8. The design rules above
 are statements about the DM axis and are unaffected, but the {\em absolute} $D$ at a subbanded
 production config is not measured; the subband ablation suggests it is $1.2$--$2.4\times$
 higher.
\item Nothing at CHORD scale has been measured. The streaming reduction removes the memory
 obstacle, but scoring $D$ with \verb|VarMapDistance| still needs the dense $A$; a variant which
 accepts $(\bar A, y, \mbox{labels})$ would close this, since for a coarse-grained-$Q$
 approximation admissibility is a theorem rather than a measurement.
\item Everything is single-tree with no early triggers, and no time-downsampled trees.
\item The shape-space greedy construction of $W$ is the one step that does not scale:
 $O(N_\beta^{2.4} N_{\rm freq})$, extrapolating to days at CHORD's $N_\beta$. Everything else is linear
 or better. It is a clustering, so a cheaper one would do, but the obvious substitutes
 (k-means on log-shape) are erratic at group granularity.
\item The {\tt Detrender2d} time half-width is a large and cheap lever on all of this: raising it
 from 2 to 8 improves $D$ at fixed rank by $2.8\times$, removes the low-DM cost concentration
 entirely, makes the SVD viable again, and costs $\sim$7\% more sweep time. Whether that
 half-width is acceptable on detrending grounds is a separate question, not addressed here.
\end{itemize}

\subsubsection{Campaign 2 setup [mostly human-authored]}
\label{ssec:campaign2_setup}

This is really interesting progress!
So far we've learned that:

\begin{enumerate}

\item Coarse-graining $A \rightarrow \bar A$ is a good approximation in practice (in the
 sense that $D(A,\mbox{lift}(\bar A))$ is small). It seems safe to assume that this should always be
 the first step, and that further steps (such as $QW$-alternation) should operate on the
 smaller $\bar A$ matrix for speed.

\item $QW$-alternation is an interesting algorithm for optimizing $D(A^{\rm true}, A^{\rm approx})$, 
 given a rank-$K$ initial guess for $A^{\rm approx}$. As currently formulated, $QW$-alternation
 doesn't change the rank $K$.

\item Truncating the SVD decomposition doesn't work on its own. (But I'd like to explore
 using the SVD decomposition in combination with other ideas -- see below.)

\end{enumerate}

\par\noindent
In campaign 2, I'd like to explore the following directions:

\begin{itemize}

\item
I realized that I never specified targets for the distance $D$ and the rank $K$, so
the problem we're exploring isn't really well-posed.
For now, let's consider rank $K=8,16,32,64,128$.
For each of these ranks (and for each choice of dedispersion+detrender config),
please explore different algorithms and try to find the smallest $D$ that you can.
This will give us a sense of the ``Pareto frontier'' of distance versus rank.

\item While I was discussing the campaign 1 results with an agent, it emerged that 
 the linear programs may have constraints that are not necessary.
 For example, the $Q$-step (step 2) may currently have a constraint $q_\beta \ge 0$.
 I don't think this is necessary (only the weaker constraint $W q_\beta \ge \bar A_{\beta,:}$).
 Please investigate further and decide whether constraints can be dropped.
 Does this significantly improve our results ($D$-values)?
 What about the $W$-step?

\item We have an analytic argument (\S\ref{sec:analytic_variance}) that $A$
 has low rank if there is no detrender. If we didn't know about this analytic
 argument, one way to ``discover'' the low rank empirically would be to compute
 the SVD singular values, and find that most of them are zero (within roundoff).

 This made me curious: what does the spectrum of singular values of $\bar A$ 
 look like empirically, if there is a detrender? Do we find that the singular
 values fall off quickly, and a low rank matrix is a good approximation (with 
 respect to the Frobenius norm -- not necessarily with respect to our distance 
 function).

\item Given a rank-$K$ approximation $A^{\rm approx} = Q W^T$, and $K' < K$,
 I'd like to explore the following idea for obtaining a rank-$K'$ approximation.
 First, SVD-decompose $A^{\rm approx} = Q_0 \Lambda_0 W_0^T$, where $\Lambda_0$ is
 $K$-by-$K$ diagonal. Then, decrease the rank to $K'$ by dropping the lowest 
 (least significant) SVD modes, obtaining $A^{\rm new} = Q_1 \Lambda_1 W_1^T$.

 As previously discussed, this may ruin the approximation, in the sense that
 $D(\bar A^{\rm true}, A^{\rm new})=\infty$. However, we can use the ``$Q$-step''
 above to update $A^{\rm new} \rightarrow Q_1' W_1$, where $Q_1'$ is a new
 matrix obtained by solving a linear program, and $W_1$ is unchanged.
 Then, $D(\bar A^{\rm true}, A^{\rm new})$ will be finite (under mild assumptions
 on $W_1$) and we can continue the $QW$-alternation algorithm from here.

 Similarly, we could also keep $Q_1$ fixed and use the ``$W$-step'' to
 update $W_1$ -- I'm not sure which variant works better.

 Note that before trying this idea, you may need to remove unnecessary
 positivity constraints from the LP (see first bullet point above), since
 SVD modes have arbitrary signs.

\item Combining ideas from the previous two bullet points, we could also try 
 initializing the low-rank decomposition $A^{\rm approx} = Q W^T$ using the
 SVD decomposition of $\bar A$ (followed by $QW$-alternation), instead of the
 clustering algorithm above (``Step 1: constructing $W$'').

 This could either be done by keeping $K$ singular values, where $K$ is the
 desired rank, or following a ``schedule'' where a larger number of singular
 values are kept, and we alternate between $QW$-alternation (to reduce the
 distance) and reducing the rank (using the idea from the previous bullet point).

\item I'm wondering whether the SVD decomposition is also useful for setting
 up the LP problems in the $Q$-step and $W$-step. For example, suppose that
 we have a decomposition $A^{\rm approx} = Q W^T$, and we want to do a $Q$-step.

 Suppose we SVD decompose $A^{\rm approx} = Q_0 \Lambda_0 W_0^T$, where $W_0$
 is semiorthogonal ($W_0^T W_0 = 1$), and then solve an LP to obtain 
 $A^{\rm new} = Q_1 W_0^T$. Is this more numerically stable, because we
 have orthogonalized columns of $W_0$ to avoid near-cancellations? Are
 there other technical advantages, such as a more principled way to set
 numerical thresholds in the LP? Please investigate.
 
 Another idea: if we find empirically that some singular values in $\Lambda_0$ 
 are very small, then it may make more sense to drop these modes (decreasing the rank),
 since the extra modes aren't ``helping'' the approximation.

\item Another idea for initialization (in addition to the clustering algorithm or
 the SVD decomposition). Let $\bar A_0$ be the coarse-grained variance map with the
 detrender removed (but the dedispersion config unchanged). Using the algorithm in
 \S\ref{sec:analytic_variance}, we have an exact low-rank decomposition $\bar A_0 = Q_0 W_0^T$.
 Suppose that we fix $W_0$ and use a $Q$-step to approximate $\bar A \approx Q_1 W_0^T$.
 Is this a useful starting point for obtaining a low-rank approximation?

 More generally, if you have any other ``natural'' ideas for generating a $W$-matrix
 (even if there is no obvious way to generate $Q$ beyond solving an LP), this could
 be interesting to investigate.

\item A small idea that may help speed up the LP solver.
 Currently, we solve with a subset of the constraints for speed, then re-evaluate
 all constraints. If any constraints fail, then we add constraints and re-solve.

 Suppose that when we evaluate the constraints, we find that they only fail by
 a small amount (say $A^{\rm approx} \ge (1-\epsilon) A^{\rm true}$ everywhere).
 Then we could solve the constraints by simply rescaling the solution by $(1-\epsilon)^{-1}$ --
 the only thing we lose is that the distance is slightly suboptimal.
 In early stages of an iterative algorithm (either $QW$-iteration or a larger iterative
 algorithm involving SVD decomposition and rank reduction) this may be a win -- it may
 be better to trade strict optimality for speed in the early stages.

\item So far we haven't analyzed dedispersion configs with \verb|subband_counts|.
Do our results (1 and 2 above) hold up in this case?
Does the distance-rank ``frontier'' get worse in the presence of \verb|subband_counts|?

\item Similarly, how do our results depend on $W$?
So far, we've only analyzed $W=2$.
Note that $W=0$ implies $n=0$.

\item Please be creative and pursue your own ideas, or extend/generalize the ideas here.
 As you explore the problem, new ideas may suggest themselves.
 Sometimes an experiment is useful if it improves general understanding of the problem domain,
 even if it doesn't immediately lead to a concrete improvement in the algorithm.

\end{itemize}

\par\noindent
In addition, here are some directions to explore that were suggested by an agent:

\begin{itemize}

\item {\bf The {\tt Detrender2d} knot vector as a design variable.}
 Campaign 1 found that the detrender does not merely perturb the structure of $A$, it
 destroys it: with no detrender, 33\% of a subbanded $A$ is exactly zero (a multiplet gets
 no contribution from channels outside its subband), and adding a single-zone detrender
 takes that to 0.0\%. But a detrender whose spline zones are {\em aligned with the frequency
 zones} keeps 25\% of the zeros. So the sparsity that a low-rank construction would exploit
 is partly a {\em choice} (the knot vector), not a given.

 This suggests a direction that changes the problem rather than the algorithm: are there
 knot vectors which leave $A$ substantially more approximable at fixed rank, and what do
 they cost in detrending quality? The time half-width $W$ is one detrender knob and is
 already covered above; the knot vector is a second, independent one, and nothing has
 explored it. If the cost is small this may be the cheapest available win.

\item {\bf Rank reduction guided by $D$ rather than by singular value.}
 This is the natural competitor to the SVD-truncation idea above, and the two should be
 measured head to head in the same experiment. Instead of dropping the modes with the
 smallest singular values, drop the column of $W$ whose removal -- followed by a $Q$-step --
 increases $D$ least, repeating until the rank is $K'$ (a ``leave-one-out'' criterion).

 The motivation is that a singular value measures a mode's contribution in Frobenius norm,
 which is the wrong norm here: campaign 1 found that pruning $W$ by column {\em
 usage} was $2.1\times$ worse than simply rebuilding at the target rank, and diagnosed it as
 keeping ``columns that supply mass, not columns that supply shapes nothing else covers''.
 A singular value is also a mass-like criterion. A leave-one-out score is aligned with the
 objective by construction. It costs $K$ $Q$-steps per removal, but those are independent
 and each is cheap.

\item {\bf Pivoted QR (or CUR) for choosing the columns of $W$.}
 Campaign 1 tried using actual rows of $A$ as columns of $W$, chosen worst-covered-first,
 and found them $1.2$--$1.4\times$ worse than the greedy clustering. That may say more about
 the selection rule than about the idea: a rank-revealing QR with column pivoting applied to
 $\bar A^T$ chooses a well-conditioned subset of rows, which is the principled version of
 the same construction.

 It has several practical advantages worth testing: it is deterministic, it costs
 $O(K N_\beta N_{\rm freq})$ rather than the $O(K_0^{2.4})$ of the agglomerative greedy, the
 atoms stay nonnegative and interpretable (each is a real output's frequency response), and
 the pivot order gives a natural nested sequence of $W$-matrices -- which is exactly what a
 rank schedule wants.

\item {\bf Warm-starting the linear programs.}
 The $Q$-step solves $N_\beta$ independent LPs, and the whole thing is repeated across
 alternations and again at every rank in a schedule. Consecutive solves differ only
 slightly, so a dual-simplex warm start from the previous optimal basis should be a large
 constant factor. Neighbouring groups are also similar to each other (adjacent coarse DMs
 have similar rows), so one group's basis can seed the next.

 This is complementary to the constraint-generation idea above rather than an alternative to
 it, and it may well be the larger of the two effects. Worth measuring as LP wall clock per
 $Q$-step, with and without.

\item {\bf Does a $W$-matrix transfer between trees and between configs?}
 Building $W$ is the one step of the recommended algorithm that does not scale:
 $O(K_0^{2.4} N_{\rm freq})$, which extrapolates to days at CHORD's $K_0$, while everything
 else is linear or better. But a {\tt DedispersionPlan} has several trees, and production
 needs weights for all of them.

 So: does a $W$-matrix built for one tree work for another (same $N_{\rm freq}$, different
 rank or subband structure)? If it does, the expensive step amortizes over trees, and
 possibly over configs. There is some reason for optimism -- campaign 1 found $W$
 need only be {\em reasonable} rather than optimal (a cheap k-means $W$ matched the
 greedy at $K=32$). Note the related question of transfer across {\em granularity} within one
 tree is already settled: building $W$ from $\bar A$ at the same $L$ used for
 $Q$ is within a few percent of building it from a finer $\bar A$, so there is no
 reason to keep a finer one around.

\end{itemize}

\subsubsection{Campaign 2 results [agent-authored]}
\label{ssec:campaign2_results}

Campaign 2 ran fourteen experiments against the questions above. The short version is that the
algorithm of \S\ref{sssec:varmap_proposed} was replaced in two places, is now $2$--$3\times$
better at every production rank, and has stopped being limited by rank at all. Two of the
bullet-point directions turned out to be dead ends and are reported as such.

\paragraph{The recipe, as it now stands.}
Three of the four steps of \S\ref{sssec:varmap_proposed} survive unchanged: coarse-grain
$A \rightarrow \bar A$ by a max over structured groups, solve one covering LP per group for
$q_\beta$, and alternate with the $W$-step. What changed is how $W$ is built, and the
{\em constraint set}:

\begin{enumerate}
\item {\bf $W$ is now a truncated SVD of $\bar A$}, not an agglomerative clustering.
 Use the SVD of the unit-sum {\em shape} matrix $S_{\beta F} = \bar A_{\beta F}/\sum_F \bar A_{\beta F}$
 for $K \gtrsim 32$ and of $\bar A$ itself below that. Only at $K \le 8$ does the clustering
 still win.
\item {\bf The $Q$-step constraint $q_\beta \ge 0$ is dropped.} It was never implied by the
 problem -- admissibility constrains the {\em product} $W q_\beta \ge \bar A_{\beta,:}$, not the
 coefficients -- and dropping it is worth $1.4$--$3.9\times$ in $D$ at fixed rank, like for like.
 The $W$-step's $W \ge 0$ may be kept when the initial $W$ is nonnegative (it is inactive at the
 optimum, worth $0$--$2.6\%$), but must {\em not} be imposed on top of a signed initializer,
 where it costs $8.8\times$.
\end{enumerate}

\par\noindent
The bullet points above asked whether the SVD could be used for initialization, for rank
reduction, and for conditioning. The answers are yes, yes, and no -- but the reason the first
answer is yes turns on a distinction the campaign took a long time to draw, and which is worth
stating plainly because it is the single most useful idea to come out of this work:

\begin{quote}
{\bf The SVD as an {\em approximation} and the SVD as a {\em $W$-matrix} are different objects.}
A truncated SVD used directly as $A^{\rm approx}$ has an admissibility cliff: below $K \approx 32$
it puts a non-positive value on a positive entry of $\bar A$, which no rescaling can repair, so
$D = \infty$. The same truncation used as the {\em $W$-matrix}, with the LP free to choose
$q$, has no cliff at all.
\end{quote}

\par\noindent
These two are easy to conflate, and conflating them is why the winning algorithm was nearly
abandoned: measurements of the first were repeatedly read as verdicts on the second. Note also
that the two changes are coupled. A signed $W$ is unusable while $q_\beta \ge 0$ is
imposed -- the sign convention of an SVD mode is arbitrary, and with $q$ constrained the LP was
infeasible on $400$ of $400$ sampled rows. With $q$ free the problem is exactly invariant under
$W_{:,c} \rightarrow -W_{:,c}$, $q_c \rightarrow -q_c$, so the sign convention stops mattering
(measured: identical $D_0$ to 17 digits). Dropping the constraint is what made an SVD $W$-matrix
available at all.

\paragraph{The frontier.}
At the real CHIME tree-0 geometry ($r=15$, $N_{\rm freq} = 16384$, coarse-graining $L_0 = 6$),
with the recipe above:

\begin{center}
\small
\begin{tabular}{lccccc}
\toprule
rank $K$ & 8 & 16 & 32 & 64 & 128 \\
\midrule
$D$, campaign 1        & 0.0412 & 0.0277  & 0.0191  & 0.0136  & 0.0105  \\
$D$, campaign 2        & 0.0376 & 0.0147  & 0.00686 & 0.00533 & 0.00462 \\
improvement            & 1.10   & 1.88    & 2.78    & 2.56    & 2.28    \\
$D/D_{\rm floor}$      & 8.52   & 3.34    & 1.55    & 1.21    & {\bf 1.05} \\
\bottomrule
\end{tabular}
\end{center}

\par\noindent
where $D_{\rm floor} = 0.00441$ is the coarse-graining floor: the covering constraint forces
$y^{\rm approx}_\alpha \ge \sum_F \bar A_{\beta F}$ for every $\alpha \in \beta$, so {\em no}
approximation with a coarse-grained $Q$, of {\em any} rank, can beat the unrescaled max-envelope over the
$L_0$ groups. The last row is the important one.

\paragraph{Rank is spent; the coarse-graining is the remaining lever.}
At $K = 128$ we are $4.8\%$ above the floor, so further rank buys almost nothing. This was
initially misread. The frontier's log-log slope softens steadily with $K$ (from $-0.57$ to
$-0.22$), which looked like an unexplained bend, and the naive diagnostic -- ``$D$ is still
$2.4\times$ the floor, so the floor is not binding'' -- is simply wrong reasoning. {\bf The floor
enters additively}, and therefore governs the slope long before $D$ approaches it. Subtracting it
leaves a clean power law,
\begin{equation}
D(K) \;\simeq\; D_{\rm floor} + c\,K^{-2}
\label{eq:varmap_c2_powerlaw}
\end{equation}
which fits over a factor 16 in rank and predicted $D(K{=}128)$ to $0.5\%$ before it was measured.
The right measure of ``how floor-limited is this cell'' is $1 - D_{\rm floor}/D$, not
$D/D_{\rm floor}$.

Since $D_{\rm floor}$ is set by $L_0$, the lever is $L_0$. Refining it from $6$ to $4$ buys
$2.10\times$ at $K = 64$, against at most $1.21\times$ available from {\em any} rank at $L_0 = 6$
-- about $7\times$ more per unit of storage. The crossover is sharp and production has just
crossed it: at $2\times$ storage from $(L_0, K) = (6, 32)$, buying rank wins $1.45\times$ to
$1.26\times$; from $(6, 64)$, buying groups wins $2.10\times$ to $1.15\times$. Below
$K \approx 29$ buying groups is a $3$--$7\times$ mistake. {\bf So: refine $L_0$ first, then buy
rank} -- and the campaign was right to optimize rank until the new $W$-matrix landed.

The axis is short, and the limit is in the shipped code rather than in the mathematics:
{\tt DedispersionPlan} accepts {\tt wt\_dm\_downsampling} $= 64, 32, 16$ but fails an assertion at
$8$, because {\tt dm\_downsampling} defaults to $2^R$. Hence $L_0 \ge R$, giving exactly two
refinements; after them $D/D_{\rm floor}$ is back to $2.3$ and rank binds again at $K \approx 75$.

\paragraph{Cost: the one step that did not scale is gone.}
The clustering construction of $W$ was $O(N_\beta^{2.4} N_{\rm freq})$ and extrapolated to days at
CHORD. The SVD $W$ is one Gram eigendecomposition: $34$~s against $\sim 2.5$~h at subscale, and
$120$--$730$~s against $\sim 44$~h at the production point. Everything else was already linear or
better. Separately, $W$ turns out not to need {\em fitting} at all: one $W$-step takes
a $W$-matrix built for a {\em different} tree, subband count or detrender to within $3$--$9\%$ of
a purpose-built one, and the analytic no-detrender factorization of \S\ref{sec:analytic_variance}
costs only $1.0$--$1.8\times$ with no clustering whatever. A {\em random} $W$ costs
$1.2$--$2.0\times$, so the choice still matters -- it just needs no expensive search.

\paragraph{Dependence on the configuration.}
\begin{itemize}
\item {\bf Subbands are expensive}: $D(N)/D(N{=}1)$ is $3.1$--$3.8\times$ at $N = 9$ and
 $5.7$--$22.8\times$ at $N = 16$--$20$, far worse than campaign 1's $1.2$--$2.4\times$ ablation
 estimate. At the CHORD-shaped point ($r = 15$, $N = 25$, 400 channels) the recipe reaches
 $D = 0.047, 0.028, 0.014$ at $K = 32, 64, 128$.
\item {\bf The detrender half-width $W$ is a real lever and does not saturate at 8}: $W = 8$ buys
 $2.5\times$ at $K = 64$ over $W = 2$, and $W = 16$ a further $1.7\times$. Whether such a window
 is acceptable on detrending grounds is still a separate question.
\item {\bf The knot vector is a weak lever, but not a dead one.} Aligning spline zones with
 {\em frequency} zones is worth nothing. Aligning them with {\em subband} boundaries is worth
 $1.1\times$ at $K = 64$ and $2.5\times$ at $K = 128$. What that alignment costs in detrending
 quality is not answerable from $A$ alone.
\end{itemize}

\paragraph{Two methodological warnings, both learned the hard way.}
These cost the campaign more than any single measurement, and both should be applied to anything
in this section:

\begin{enumerate}
\item {\bf A conclusion measured at $N_{\rm freq} = 400$ may reverse at $16384$.} It happened
 three times: the SVD-versus-LP comparison, the subband rank-budgeting rule, and ``the $W$-step is
 redundant''. The mechanism is understood -- the LP family gets $\sim 3\times$ {\em better} with
 more channels while the SVD family gains nothing -- but the practical rule is simply to treat any
 subscale recommendation as provisional.
\item {\bf A floor-based conclusion is not an LP-based conclusion.} The coarse-graining floor
 ranks the $W$ ladder perfectly and the knot-vector axis {\em oppositely}: one configuration with
 a $19\times$ better floor is $1.3\times$ {\em worse} in achieved $D$. Use the floor as a bound,
 never as a ranking. The singular-value spectrum of $\bar A$, by contrast, {\em does} rank
 configurations correctly (Spearman $\rho \approx 0.97$--$0.99$) for the cost of one
 eigendecomposition, and is a good cheap triage -- though it too degrades above $K \approx 64$,
 where the floor takes over.
\end{enumerate}

\paragraph{What is still open.}
The largest item is not a physics question but a solver one. The covering LPs are solved by HiGHS,
and at high rank on heavily subbanded maps a small fraction of them ($\lesssim 1\%$) return a
point that violates their own constraints while reporting success, or fall back to a crude seed.
This is not a rounding-level nuisance: one failure per $\sim450$ groups costs more $D$ than an
entire doubling of the rank, and it is why the refined-$L_0$, $K = 128$ cells are currently
estimates rather than measurements. A repair is known (halve the rank, re-solve, keep the rank)
and takes seconds; until it is in place, any number at $K \ge 128$ on a subbanded map should be
treated as provisional. Beyond that: warm-starting the LPs was never attempted, and nothing has
been measured at CHORD's channel count.

\subsubsection{Campaign 3 setup [human + agent authored]}
\label{ssec:campaign3_setup}

Directions for the next campaign. (The rest of the campaign-2 leftovers -- the solver defect at
$K \ge 128$, the exact $(L_0, K) = (4, 128)$ cell, LP warm starts, and CHORD's channel count --
are tracked outside these notes.)

The campaign 2 results look promising!
The main things I'm worried about are {\bf scalability} and {\bf robustness}.
Consider the CHORD scale config \verb|chord_sb2.yml| (primary tree 0 only), with
$N_{\rm freq} = 28160$, rank $r=15$, and 25 subbands.
Can we reach this scale in a reasonable amount of computing time (say $<1$ day on 32 cores)?
Will our current algorithm run into numerical instability (e.g.\ in the LP solver) before
reaching this scale?
What rank-$D$ ``Pareto frontier'' will we see at CHORD scale?

\begin{itemize}

\item Instead of running \verb|chord_sb2.yml| directly, approach it via a sequence of
 increasingly large-scale configs (e.g.\ increasing $r$ by 1 and/or doubling $N_{\rm freq}$,
 and/or adding subbands. You may find that new issues emerge as the scale increases.

\item Try to reach as large a scale as possible, within the time horizon of the ch-research
 skill. In order to achieve this goal, consider carefully what set of parallel subagents will
 make best use of our 32 cores and 2 GPUs. For example, if you're precomputing $A$-matrices
 on the GPU and solving LPs on the CPU, you may want to precompute the next problem instance
 in parallel with solving the current instance.

\item You may not reach CHORD scale within the time horizon. In that case, please try to
 find an empricial scaling law based on your subscale results, and predict what ``Pareto
 frontier'' we might see at CHORD scale.

\item This idea is cut-and-paste from an earlier section -- I'm including it again since
 it may be more helpful at scale.

 A small idea that may help speed up the LP solver.
 Currently, we solve with a subset of the constraints for speed, then re-evaluate
 all constraints. If any constraints fail, then we add constraints and re-solve.

 Suppose that when we evaluate the constraints, we find that they only fail by
 a small amount (say $A^{\rm approx} \ge (1-\epsilon) A^{\rm true}$ everywhere).
 Then we could solve the constraints by simply rescaling the solution by $(1-\epsilon)^{-1}$ --
 the only thing we lose is that the distance is slightly suboptimal.
 In early stages of an iterative algorithm (either $QW$-iteration or a larger iterative
 algorithm involving SVD decomposition and rank reduction) this may be a win -- it may
 be better to trade strict optimality for speed in the early stages.

\item {\bf Pin a nonnegative column in $W$, and exclude it from the $W$-step.}
 Here's a minor suggestion that may be useful.

 Campaign 2 dropped $q_\beta \ge 0$ entirely and, under the SVD recipe, dropped $W \ge 0$ as
 well, so both factors are now sign-free. That is where the performance came from, but it
 quietly removed something several parts of the algorithm still rely on: {\em the existence of
 at least one all-positive column of $W$}. Three places currently hope for one rather than
 requiring it -- the additive admissibility repair (which raises the product in every channel at
 once only if the column it adds is nonnegative), the one-hot seed for the first $Q$-step, and
 the argument that a covering LP is feasible at all. Today that hope is usually satisfied only
 because the leading SVD mode of a nonnegative matrix is itself nonnegative (Perron--Frobenius),
 which is a fact about {\em one} mode and is not guaranteed to survive truncation, shape
 normalization, or a differently-constructed $W$.

 The proposal is to stop relying on it: reserve one column of $W$, fix it to a nonnegative
 vector, and exclude it from the $W$-step so it cannot drift. The cost is bounded by one unit of
 rank, and is probably less -- see below.

 {\em Which vector.} An all-ones column is the obvious choice and is the wrong one. It is
 nonnegative and dominating, but its cost in the objective is $s_c = N_{\rm freq}$, so covering
 a group whose largest entry is $m$ costs $m N_{\rm freq}$ against a true row sum of
 $\sum_F \bar A_{\beta F}$ -- i.e.\ one pays the row's peak-to-mean ratio. On a subbanded map a
 row is supported on roughly $1/N$ of the band, so at CHORD's $N = 25$ the all-ones column is
 $\sim 25\times$ more expensive than the row it is covering, and the LP would never choose it
 except in extremis. A better choice is the elementwise column max over groups,
 \begin{equation}
 W_{F,c_0} = \max_\beta \bar A_{\beta F}
 \label{eq:varmap_pinned_column}
 \end{equation}
 which is nonnegative by construction, is positive in exactly the channels where some group
 needs it, costs one streaming pass to build (the same pass that already produces $\bar A$), and
 is the {\em tightest} single vector that dominates every group. It is also the rank-1 member of
 campaign 1's envelope family, so unlike an all-ones column it is a vector the LP will actually
 use on hard groups.

 {\em What it guarantees.} With Eq.\ (\ref{eq:varmap_pinned_column}) pinned, $q_\beta = e_{c_0}$
 is a certified feasible point for {\em every} group, with a nonnegative product, whatever the
 other $K-1$ columns look like. We currently have no such certificate, which is why campaign 2
 had to prove feasibility by hand to establish that the solver was reporting a feasible LP as
 infeasible. It also means the one-hot seed can never fail to exist, and the additive repair can
 never run out of usable columns.

 {\em What it does not fix.} It does not prevent the product $W q_\beta$ from going negative in
 channels where $\bar A_{\beta F} = 0$ -- that is what constraining every channel is for -- it
 only makes such a case repairable. And it does not address the solver's rank-dependent failures
 at $K \ge 128$, which are a conditioning problem, not a feasibility one. It is complementary to
 both: a certified feasible point is exactly what one would hand the solver as a warm start.

 {\em The measurement.} Compare the current recipe at rank $K$ against the pinned variant at the
 same rank $K$ (so $K-1$ learned columns plus one pinned), on at least one subbanded and one
 unsubbanded map, reporting $D$ and the LP failure counts. The interesting question is whether
 the pinned column costs a full unit of rank or a fraction of one: the leading SVD mode is
 already approximately the mean spectrum, so a pinned envelope may be largely reconstructing
 something the truncation would have spent a mode on anyway.

\item Does it help (e.g.\ with numerical instability at scale) to ``re-orthogonalize'' the
 rank-$K$ approximation after each $Q$-step or $W$-step, by taking the SVD decomposition? 
 (Note: if we decide to ``pin'' one or more all-positive columns, as suggested in the previous
 bullet point, then we'd want to exclude pinned columns from the re-orthogonalization.

\item Feel free to pursue any loose ends from campaigns 1 and 2 that weren't fully resolved,
 or that you want to revisit. 
 In particular, ideas that were abandoned previously in subscale instances may prove
 to be useful at larger scale.

\item
As usual, please get creative and explore your own ideas!

\end{itemize}

\subsection{Brute-force computation of the variance map}
\label{ssec:variance_map_bruteforce}

Although Eq.\ (\ref{eq:variance_map_def}) defines $A$ through a random ensemble, $A$ is really just
a statement about the matrix elements of $L$, and can be computed one column at a time by applying
$L$ to ``one-hot'' input arrays. This subsection works out that brute-force algorithm.

It is very slow, but it treats every stage of $L$ as a black box, so it needs no new analysis to
handle subbands, time-downsampled trees, early triggers, or -- most importantly -- the
{\tt Detrender2d}, which the machinery of \S\ref{sec:analytic_variance} cannot currently handle at
all. In fact the relationship is closer than that: \S\ref{sec:analytic_variance} {\em is} this
algorithm, executed sparsely (propagating the one-hot impulse response through the dedispersion
steps in compressed form) instead of densely. The brute-force version is the dense reference that
the sparse version should be tested against.

Besides the parameters of \S\ref{ssec:variance_map_setup}, we will need the {\tt Detrender2d} time
half-width $W$ (\dtnotes{2-d detrending}) and the max peak-finding kernel width $W_{\rm max}$
(\ddnotes{Peak-finding}).

\subsubsection{$A$ as a set of column norms}

Write the matrix elements of $L$ as $L_{\alpha t, F t'}$, i.e.
\begin{equation}
{\tt out}_\alpha[t] = \sum_{F=0}^{N_{\rm freq}-1} \sum_{t'} L_{\alpha t, F t'} \, {\tt in}[F,t']
\label{eq:bf_L_matrix}
\end{equation}
where $0 \le \alpha < 2^{r-R}MP$ is the flattened output index, $t$ is a downsampled output time,
and $t'$ is an undownsampled input time. Since the input has covariance
$\langle {\tt in}[F,t'] \, {\tt in}[F',t''] \rangle = v_F \, \delta_{FF'} \, \delta_{t't''}$,
\begin{equation}
y_\alpha = \big\langle {\tt out}_\alpha[t]^2 \big\rangle = \sum_F v_F \sum_{t'} L_{\alpha t, F t'}^2
\end{equation}
and comparing with Eq.\ (\ref{eq:variance_map_def}),
\begin{equation}
A_{\alpha F} = \sum_{t'} L_{\alpha t, F t'}^2
\label{eq:bf_A_colnorm}
\end{equation}
for any output time $t$ which has reached steady state. That is, {\bf a column of $A$ is a set of
squared column norms of $L$}, taken over the input time axis at fixed output time.

Now let $e^{(F,t_0)}$ denote the one-hot input array ${\tt in}[F',t'] = \delta_{FF'} \delta_{t't_0}$.
Applying $L$ to it picks out a {\em row} of $L$, not a column:
\begin{equation}
\big( L \, e^{(F,t_0)} \big)_\alpha[t] = L_{\alpha t, F t_0}
\hspace{1.2cm}\Longrightarrow\hspace{1.2cm}
\sum_t \big( L \, e^{(F,t_0)} \big)_\alpha[t]^2 = \sum_t L_{\alpha t, F t_0}^2
\label{eq:bf_rownorm}
\end{equation}
So the proposed algorithm computes a {\em row} norm where Eq.\ (\ref{eq:bf_A_colnorm}) asks for a
{\em column} norm. The two agree by time-translation invariance -- but the details differ between
$\gamma=0$ and $\gamma>0$, which is the one place where the algorithm needs care.

\subsubsection{Undownsampled trees ($\gamma=0$)}

Every stage of $L$ is time-translation invariant. The {\tt Detrender2d} run with no mask is a fixed
FIR filter along time: with no mask its normal equations, and hence its mask expansion and its
conditioning cuts, depend on neither the data nor $t$, so $L$ really is both linear and
shift-invariant. Tree gridding acts only on the frequency axis, and tree dedispersion and the
peak-finding convolutions are sums of lagged copies. Therefore
\begin{equation}
L_{\alpha t, F t'} = \ell_{\alpha F}[t-t']
\end{equation}
and Eqs.\ (\ref{eq:bf_A_colnorm}), (\ref{eq:bf_rownorm}) both collapse to
$\sum_u \ell_{\alpha F}[u]^2$. Hence, for any $t_0$ satisfying the support conditions of
\S\ref{sssec:bf_sizing} below,
\begin{equation}
A_{\alpha F} = \sum_t \big( L \, e^{(F,t_0)} \big)_\alpha[t]^2
\label{eq:bf_master}
\end{equation}
i.e.\ apply $L$ to a one-hot input, square, and sum over the output time axis. One pass per input
channel gives all $N_{\rm freq}$ columns of $A$.

\subsubsection{Downsampled trees ($\gamma>0$)}

For $\gamma>0$, the operator $L$ contains a time-downsampling stage $\TDS(\gamma)$
(\ddnotes{Downsampled trees and the LaggedDownsampler}) between the detrender and the tree, and is therefore invariant
only under input shifts by $2^\gamma$ samples. The right generalization is a {\em polyphase}
decomposition of the input time index into a downsampled index and a phase,
\begin{equation}
t' = 2^\gamma s + c, \hspace{1.5cm} 0 \le c < 2^\gamma
\end{equation}
Invariance under $s \rightarrow s+1$ (which shifts the output by one downsampled sample) gives
\begin{equation}
L_{\alpha t, F(2^\gamma s + c)} = \ell^{(c)}_{\alpha F}[t-s]
\end{equation}
with $2^\gamma$ generally-distinct impulse responses $\ell^{(0)},\cdots,\ell^{(2^\gamma-1)}$.
Eq.\ (\ref{eq:bf_A_colnorm}) now reads
\begin{equation}
A_{\alpha F} = \sum_{c=0}^{2^\gamma-1} \sum_s \ell^{(c)}_{\alpha F}[t-s]^2
  = \sum_{c=0}^{2^\gamma-1} \big\| \ell^{(c)}_{\alpha F} \big\|^2
\end{equation}
whereas a single one-hot at input time $t_0$ sees only the phase $c = t_0 \bmod 2^\gamma$.
So a downsampled tree does need $2^\gamma$ one-hot arrays per input channel, one per residue class
of $t_0$ modulo $2^\gamma$, and the results are {\em summed} (not averaged):
\begin{equation}
A_{\alpha F} = \sum_{c=0}^{2^\gamma-1} \sum_t \big( L \, e^{(F,t_c)} \big)_\alpha[t]^2,
  \hspace{1.5cm} t_c \equiv c \ (\mbox{mod } 2^\gamma)
\label{eq:bf_master_ds}
\end{equation}
Any representative $t_c$ of each residue class will do; e.g.\ $t_c = t_0 + c$ with $t_0$ a multiple
of $2^\gamma$.

There is one important special case. If the part of $L$ upstream of $\TDS(\gamma)$ is
{\em instantaneous} in time -- no detrender at all, or a {\tt Detrender2d} with $W=0$, which fits
each time sample independently and is therefore a pure frequency-axis operator -- then
$\TDS(\gamma)$ maps $e^{(F, 2^\gamma s + c)}$ to the {\em same} downsampled array for every $c$,
the $2^\gamma$ phases coincide, and one pass suffices:
\begin{equation}
A_{\alpha F} = 2^\gamma \sum_t \big( L \, e^{(F,t_0)} \big)_\alpha[t]^2
  \hspace{1.5cm} (W=0)
\end{equation}
This is the familiar statement that downsampling with a $2^{-\gamma/2}$ normalization leaves the
variance unchanged (\ddnotes{Time-downsampled trees and early triggers}), and it is worth keeping as a unit test: with
$W=0$ the $2^\gamma$ phase outputs must agree {\em exactly}. For $W>0$ they genuinely differ, since
the polyphase components of the detrender's time response differ, and all $2^\gamma$ passes are
needed.

\subsubsection{Placing the one-hot, and sizing $N_{\rm time}$}
\label{sssec:bf_sizing}

Eqs.\ (\ref{eq:bf_master}), (\ref{eq:bf_master_ds}) are exact only if the {\em entire} impulse
response is emitted: anything clipped at either end of the array is silently dropped from the sum
of squares, and the result is an underestimate with no error indication. Three conditions.

\begin{itemize}

\item {\bf Detrender padding.} The {\tt Detrender2d} is handed a length-$(T+2W)$ buffer and
overwrites only samples $[W, W+T)$ with the residual (\dtnotes{2-d detrending}). A one-hot at buffer
index $t_0$ produces a detrended response on $[t_0-W,\, t_0+W]$, all of which must lie in the
emitted region, so we need $t_0 \ge 2W$.

\item {\bf No negative times.} Tree dedispersion defines ${\tt in}[\cdots,t] = 0$ for $t<0$
(\ddnotes{Tree dedispersion}), so a response reaching back before the start of the stream would
be truncated -- equivalently, we would be computing a non-steady-state variance
(cf.\ the footnote to Eq.\ \ref{eq:avar_master}). The previous condition already rules this out,
since the response starts at $t_0 - W \ge W \ge 0$.

\item {\bf Enough room downstream.} Going forward from the one-hot, the response is spread by
three stages:
\begin{itemize}
\item the detrender, over $2W+1$ undownsampled samples, hence
$\lfloor 2W/2^\gamma \rfloor + 1$ downsampled samples after $\TDS(\gamma)$;
\item (subbanded) tree dedispersion, by at most $\Delta_{\rm dd}$ downsampled samples, where
\begin{equation}
\Delta_{\rm dd} \equiv \begin{cases} 2^r & \mbox{if } \gamma = 0 \\
                                     2^{r+1} & \mbox{if } \gamma > 0 \end{cases}
\label{eq:bf_ddspread}
\end{equation}
\item the peak-finding convolution, by at most $T_{\rm max} - 1 = 2W_{\rm max}-1$ downsampled
samples (\S\ref{sec:avar_convolver}). If the kernel applies a per-profile time shift
(\ddnotes{Dedispersion output arrays}), that only translates the response, which is harmless
here.
\end{itemize}
\end{itemize}
\par\noindent
Eq.\ (\ref{eq:bf_ddspread}) deserves a comment. $\Delta_{\rm dd}$ is simply the largest full-band
delay searched by the tree, because the output time index is by construction the arrival time
extrapolated to the top of the tree's own band (\ddnotes{Subband search}) -- so the per-subband
time lags $T$ do not {\em add} to this bound, they are part of it. Concretely, in Case 1 of
subbanded dedispersion (\ddnotes{Subband search}) the total lag from tree-freq $f$ to the
output is at most
\begin{equation}
\underbrace{T}_{= \, d_{\rm hi}2^l(2^{R-l}-1-f)} + \underbrace{d}_{= \, d_{\rm hi}2^l + d_{\rm lo}}
\ \le \ d_{\rm hi} 2^R + d_{\rm lo} \ \le \ (2^{r-R}-1)2^R + (2^l-1) \ < \ 2^r
\end{equation}
(maximized at $f=0$, and using the fact that the largest within-subband lag is
$L_{r-R+l}(0,d) = d$), and Case 2 is the same. The two cases in Eq.\ (\ref{eq:bf_ddspread}) differ
because for $\gamma>0$ the tree keeps only the upper DM half (\ddnotes{Time-downsampled trees and early triggers}), so
the tree rank $r$ is one less than the rank of the underlying $\DD$: in the code,
$r = {\tt total\_rank} = r_{\rm top} - e - [\gamma>0]$, and $\Delta_{\rm dd} = 2^{r_{\rm top}-e}$
downsampled samples in both cases.

Adding the three contributions, the response occupies at most
\begin{equation}
n_t = \left\lfloor \frac{2W}{2^\gamma} \right\rfloor + \Delta_{\rm dd} + 2 W_{\rm max}
\end{equation}
downsampled output samples starting from the one-hot, so it suffices to take
\begin{equation}
2W \le t_0 < 2W + 2^\gamma
  \hspace{1.5cm}
N_{\rm time} \ \ge \ 4W + 2^\gamma \big( \Delta_{\rm dd} + 2W_{\rm max} + 1 \big)
\label{eq:bf_ntime}
\end{equation}
where the $4W$ is $2W$ from the placement of $t_0$, plus $W$ for the detrender's forward reach,
plus $W$ of trailing detrender padding; and the ``$+1$'' inside the parentheses absorbs the choice
of phase $c$.

Two practical remarks. First, because the answer is a sum over {\em all} output times, the
algorithm is insensitive to how the stream is chunked: one may equally run the dedisperser over
several chunks totalling $N_{\rm time}$ samples and accumulate $\sum_t(\cdots)^2$ across them, and
there is no need to know when each DM row reaches steady state. Second, the cheapest guard against
a mis-sized array is to assert that the last few emitted output samples are identically zero;
Eq.\ (\ref{eq:bf_ntime}) is a bound, not a tight estimate, and it is easy to get wrong by a factor
of $2^\gamma$.

\subsubsection{Cost, and relation to the other algorithms}

Each pass costs one full evaluation of $L$ on an $(N_{\rm freq}, N_{\rm time})$ array, and
$N_{\rm freq} 2^\gamma$ passes are needed per tree ($N_{\rm freq}$ passes if $W=0$). For a
CHORD-sized config that is $28160$ passes per tree, each dedispersing $\sim\!10^5$ downsampled time
samples at rank $r=16$: hopelessly slow. Brute force is a small-config reference implementation,
not a production algorithm -- and not even a practical replacement for
\S\ref{sec:analytic_variance}, which is itself described there as too slow for production. Two
things make it less bad than it looks: passes for different $F$ can be batched along the
dedisperser's spectator beam axis, and almost all of the work in each pass is on identically-zero
data, which a specialized implementation could skip. Skipping it systematically is precisely what
\S\ref{sec:avar_sparsetile} does.

Three remarks on how this fits with the rest of the notes.

\begin{itemize}

\item {\bf It generalizes the Monte Carlo check} of \S\ref{sec:avar_validation}. That check feeds
random Gaussian noise with per-channel variance $V_F$ through a {\tt ReferenceDedisperser} and
compares empirical variances against the analytic prediction. The brute-force algorithm runs the
same forward machinery, but replaces the random ensemble by the deterministic standard basis
$\{e^{(F,t_0)}\}$ -- so the result is exact rather than noisy, and it yields the full matrix $A$
rather than a single product $Av$.

\item {\bf It should agree with {\tt PfAvarExact}} wherever both apply. With the detrender
disabled, $\gamma=0$, and subbands disabled, Eq.\ (\ref{eq:bf_master}) must reproduce
Eq.\ (\ref{eq:avar_master}) column by column: indeed $A_{(d,p),F} = \|h_p * G_{Fd}\|^2$ is exactly
the statement that the two agree. This is a strong end-to-end test, since the two computations
share no code.

\item {\bf It is the only algorithm we have which handles the detrender.} The efficiency of
\S\ref{sec:avar_sparsetile} rests on the gridding footprint ${\mathcal S}_F$ being a contiguous
range of a few tree-freq channels. A {\tt Detrender2d} destroys this: the detrended response to a
one-hot in channel $F$ is spread over every channel of $F$'s spline zone (\dtnotes{2-d detrending}),
so the {\tt SparseTile} representation would degrade towards a dense one. Brute force is completely
insensitive to it. Note also that these off-diagonal responses are individually small but not
negligible in aggregate: each contributes $O(\lambda_F^2)$ to $A$, but there are $O(N_{\rm freq})$
of them, so they are a percent-level effect rather than a rounding error.

\end{itemize}
\par\noindent
Finally, three checks worth asserting in any implementation. Every entry of $A$ is a sum of
squares, so $A_{\alpha F} \ge 0$ identically. For an early-triggered tree ($e>0$, and likewise for
a tree whose subbands do not span the full band), the columns of all input channels outside the
tree's frequency range must vanish identically. And the arithmetic should be done in float64: the
one-hot response has amplitude $O(2^{-r/2})$ before squaring, and the detrender's off-diagonal
contributions are smaller still, so in float32 the aggregate effect noted above would be at risk of
being lost to roundoff against the diagonal term.

\clearpage

\appendix

\section{Monotonicity of the variance map in the DM bits (no detrender)}
\label{sec:dyadic_block}

{\bf This appendix treats the no-detrender case only.} The linear operator is bare tree
gridding $\to$ subbanded tree dedispersion $\to$ peak-finding convolution, as in
\S\ref{sec:analytic_variance}, \S\ref{sec:avar_statement} -- {\em not} the detrended
operator of \S\ref{ssec:variance_map_setup}, which includes a {\tt Detrender2d}.
{\bf The result below is proved in that case only, and the restriction is not a
technicality:} with a {\tt Detrender2d} in front of the dedisperser the property is not
merely unproved but {\em false}, and badly so (\S\ref{ssec:dyadic_sharpness}).

Fix one column of the variance map: a {\tt DedispersionTree}, a frequency subband, a
peak-finding profile $p$, and a single input frequency channel $F$. What is left is a 1-d
array $A[d_m]$ indexed by the subband's dispersion measure
$d_m = 2^l d_{\rm hi} + d_{\rm lo}$ (\ddnotes{Subband search}). The purpose of this
appendix is to prove that $A$ is {\em antitone in the binary digits of $d_m$}: writing
$d \subseteq d'$ for the bitwise-subset partial order on DM indices (i.e.\ $d_b \le d'_b$
for every bit $b$, equivalently $d \,\&\, d' = d$),
\begin{equation}
d \subseteq d' \hspace{5mm} \Longrightarrow \hspace{5mm} A[d'] \;\le\; A[d].
\label{eq:bitwise_property}
\end{equation}
Equivalently: {\em setting a bit of the DM index never increases the variance},
$A[d \,|\, 2^b] \le A[d]$. Two consequences are worth stating separately. First, taking
$d = 0$, the DM index $d_m=0$ maximizes the column. Second, and more useful, the
{\em dyadic-block property}: in every aligned dyadic block of $d_m$ indices, the {\em first}
entry is the largest,
\begin{equation}
A\big[ 2^k j \big] \;=\; \max_{0 \le i < 2^k} A\big[ 2^k j + i \big]
  \qquad \mbox{for every $k \ge 0$, $j \ge 0$ with indices in range,}
\label{eq:dyadic_property}
\end{equation}
which follows from (\ref{eq:bitwise_property}) because $2^kj$ is the bitwise-minimal element
of its own dyadic block ($2^kj$ and $i < 2^k$ have disjoint bits). The practical consequence
is that an approximation to the variance map which resolves only the top DM bits, and which
must never {\em under}estimate a variance (\S\ref{ssec:distance_function}), can simply be
evaluated at the bottom of each dyadic DM block: no maximization over the block is needed,
and the resulting bound is tight.

Concretely, $A$ is a reindexing of the arrays {\tt PfAvarExact} produces: with
$m = (\mbox{subband}, d_{\rm lo})$ and $d = d_{\rm hi}$, the per-multiplet variance
$y_m[d,p]$ of \S\ref{ssec:svd_exact_map} is $A[2^l d_{\rm hi} + d_{\rm lo}]$, so a dyadic
block of $A$ generally runs across several multiplets.

Note that $A$ is {\em not} monotone in $d_m$ as an integer: it wanders up and down, but never
above the value at any bitwise-subset index (\S\ref{ssec:dyadic_sharpness}).

The proof splits into a statement about the lag function $L_r(f,d)$ (Lemmas 2--5: {\em the
lag function is linear in the bits of $d$ with non-negative coefficients, so setting a DM bit
can only spread the impulse response out}) and a statement about the peak-finding kernels
(Lemma 6: {\em spreading out can only decrease the variance}). Both are needed;
\S\ref{ssec:dyadic_sharpness} shows that neither hypothesis can be dropped, and that the
bitwise-subset order is exactly the right partial order. {\bf The first hypothesis --
non-negativity of the gridding weights -- is exactly what a detrender destroys, which is
why the no-detrender assumption above cannot be relaxed.}

\subsection{Setup}
\label{ssec:dyadic_setup}

We work in the no-detrender setting of \S\ref{sec:avar_statement}: the linear operator is
(tree gridding) $\to$ (subbanded tree dedispersion) $\to$ (convolution with the peak-finding
profile $h_p$), with no normalization to sigmas and no coarse-graining, and the dedisperser
input has uncorrelated entries with per-channel variance $V_F$.

Fix a {\tt DedispersionTree} with early-trigger level $e$ and primary-tree index $\gamma$,
and write
\begin{equation}
\hat r \;\equiv\; r_{\rm top} - e
\end{equation}
for the rank of the dedispersion it performs: the tree dedisperses the lowest $2^{\hat r}$
tree-freq channels of the band (\ddnotes{Time-downsampled trees and early triggers}). (In the code, {\tt tree\_r}
is $\hat r$ for $\gamma = 0$ and $\hat r - 1$ for $\gamma > 0$, because the upper-DM-half
restriction is folded into the rank.) Fix a subband $(l,s)$, $0 \le l \le R$, and let
\begin{equation}
K \;\equiv\; \hat r - R + l,
  \hspace{1.5cm}
a \;\equiv\; 2^{\hat r - R} \times (\mbox{index of the subband's first coarse-freq block}),
\end{equation}
so that the subband occupies the $2^K$ contiguous tree-freq channels $a \le f < a+2^K$
(\ddnotes{Subband search}) and its DM index runs over
$0 \le d_m < 2^K$. Finally let
\begin{equation}
w_u \;\equiv\; w_{(a+u)F} \;\ge\; 0
  \hspace{1.5cm} (0 \le u < 2^K)
\end{equation}
be the tree-gridding weights (\ddnotes{Tree gridding kernel}) of input channel $F$ on the subband's
channels. Non-negativity is the only property of them we use (they are overlap lengths; that
their support is a short contiguous run is not needed).

\medskip
\par\noindent
{\bf Lemma 1 (closed form).}
{\em For a $\gamma=0$ tree, the variance-map column is}
\begin{equation}
A[d_m] \;=\; 2^{-K} \sum_{u=0}^{2^K-1} \sum_{v=0}^{2^K-1}
   w_u \, w_v \; A_p\big( L_K(u,d_m) - L_K(v,d_m) \big),
\label{eq:dyadic_closedform}
\end{equation}
{\em where $L_K$ is the lag function of \ddnotes{Non-recursive formulation of tree dedispersion}
and $A_p$ is the autocorrelation of the
peak-finding kernel $h_p$. For a $\gamma>0$ tree, the exported array is the upper DM half of
(\ref{eq:dyadic_closedform}): $A[d_m] = (\mbox{RHS evaluated at } d_m + 2^{K-1})$.}

\medskip
\par\noindent
{\em Proof.} Both cases of subbanded dedispersion (\ddnotes{Subband search}) compute
exactly $\DD(K)$ applied to the subband's $2^K$ tree-freq channels, followed by a time shift.
In Case 1 ($l=0$ or even $s$) the subband output is read off the intermediate array after
$K = \hat r - R + l$ dedispersion steps, whose coarse-freq block $f$ is $\DD(K)$ applied to
tree-freq channels $2^K f \le \cdot < 2^K (f+1)$ (Observation 1 of
\ddnotes{The non-recursive formulation of tree dedispersion}). In Case 2 ($l>0$, odd $s$) the two adjacent level-$(K-1)$
blocks $2f+1$ and $2f+2$ are combined by one further $\DDS$ step, in which block $2f+1$ plays
the ``lower'' role and $2f+2$ the ``upper'' role. This pair is not aligned in the ambient
tree, but the $\DDS$ formulas do not depend on the block index (Observation 1), so
the last-step identity $\DD(K) = \DDS(K-1) \circ (1 \otimes \DD(K-1))$ applies to the pair
verbatim: the result is $\DD(K)$ applied to
their union, which is precisely the subband.

In both cases the rank-$K$ DM index of that $\DD(K)$ is the subband DM index
$d_m = 2^l d_{\rm hi} + d_{\rm lo}$. In Case 1 this is immediate ($d = d_{\rm hi}2^l +
d_{\rm lo}$ in the notation of \ddnotes{Subband search}). In Case 2 the extra $\DDS$ step
maps the level-$(K-1)$ index $d = d_{\rm hi}2^{l-1} + d_{\rm lo}$ to $2d + b$, i.e.\ to
$2^l d_{\rm hi} + (2 d_{\rm lo} + b)$, and $2d_{\rm lo}+b$ is exactly the fine index of the
multiplet which receives it.

Therefore, by the non-recursive form of $\DD(K)$ (\ddnotes{Non-recursive formulation of tree
dedispersion}) and the one-hot
impulse response (\ref{eq:avar_impulse}), the impulse response of input channel $F$ in this
subband is
\begin{equation}
G_{F,d_m}[t] \;=\; 2^{-K/2} \sum_u w_u \, \delta_{t,\, L_K(u,d_m) + T},
\label{eq:dyadic_impulse}
\end{equation}
and Eqs.\ (\ref{eq:avar_master}), (\ref{eq:avar_autocorr}) give (\ref{eq:dyadic_closedform}),
per unit input variance $V_F$. The subband extraction lag $T$ drops out because
$\|h_p * g\|^2$ is translation invariant; note that this holds for {\em any} $T$, so the
variance map is unchanged by the alternative time-lag convention discussed in the Remark of
\ddnotes{Subband search}.

For $\gamma>0$, the tree time-downsamples its input by $2^\gamma$ before dedispersing.
Downsampling adds $2^\gamma$ samples and multiplies by $2^{-\gamma/2}$, so the downsampled
array again has uncorrelated entries with per-channel variance $V_F$, and the argument above
applies verbatim on the downsampled time grid. The only other difference is that the tree
keeps just the upper half of the coarse DM range (\ddnotes{Time-downsampled trees and early triggers}): the top bit
of $d_{\rm hi}$ is fixed to 1 and dropped, which is the DM range $d_m \ge 2^{K-1}$. $\square$

\medskip
\par\noindent
Two remarks. First, early triggers have disappeared into the definition of $\hat r$: a tree
with $e>0$ performs ordinary rank-$\hat r$ dedispersion on the first $2^{\hat r}$ tree-freq
channels, i.e.\ the same statement with a truncated channel map (this is what
{\tt PfAvarExact} does; cf.\ its {\tt cmap\_rank}). Second, only $w_u \ge 0$ is used below,
so everything that follows holds for an arbitrary channel map and an arbitrary set of subband
counts $C_l$, and -- summing (\ref{eq:dyadic_closedform}) against $V_F \ge 0$ -- for any
non-negative combination of input channels.

\subsection{Setting a DM bit spreads the impulse response out}
\label{ssec:dyadic_lag}

\par\noindent
{\bf Lemma 2 (bitwise form of $L$).}
{\em Write $d_m$ for bit $m$ of $d$. Then}
\begin{equation}
L_r(f,d) \;=\; \sum_{m=0}^{r-1} \big( 1 - f_{r-1-m} \big) \, c^{(m)}(d),
  \hspace{1.2cm}
c^{(m)}(d) \;\equiv\; \Big\lfloor \frac{d}{2^{m+1}} \Big\rfloor + d_m \;\ge\; 0,
\label{eq:dyadic_bitwise}
\end{equation}
{\em and for every $0 \le m_0 < r$ (with $0 \le d < 2^r$),}
\begin{equation}
\sum_{m=m_0+1}^{r-1} c^{(m)}(d) \;=\; \Big\lfloor \frac{d}{2^{m_0+1}} \Big\rfloor
  \;=\; c^{(m_0)}(d) - d_{m_0}.
\label{eq:dyadic_telescope}
\end{equation}
{\em In particular $L_r(0,d) = d$ and $L_r(2^r-1,d) = 0$.}

\medskip
\par\noindent
{\em Proof.} In the definition of the lag function $L_r(f,d)$ (\ddnotes{Non-recursive
formulation of tree dedispersion}), the coefficient of $(1-f_i)$ is
\begin{equation}
\sum_{j=0}^{r-1} \Omega_{i+j-r} \, d_j
 \;=\; \Omega_{-1} \, d_{r-1-i} + \sum_{j \ge r-i} 2^{i+j-r} d_j
 \;=\; d_{r-1-i} + 2^{i-r} \cdot 2^{r-i} \Big\lfloor \frac{d}{2^{r-i}} \Big\rfloor
 \;=\; d_m + \Big\lfloor \frac{d}{2^{m+1}} \Big\rfloor
\end{equation}
with $m \equiv r-1-i$, which is (\ref{eq:dyadic_bitwise}). For
(\ref{eq:dyadic_telescope}), use $\lfloor d/2^m \rfloor = 2 \lfloor d/2^{m+1} \rfloor + d_m$
to rewrite $c^{(m)}(d) = \lfloor d/2^m \rfloor - \lfloor d/2^{m+1} \rfloor$; the sum
telescopes to $\lfloor d/2^{m_0+1} \rfloor - \lfloor d/2^r \rfloor$, and
$\lfloor d/2^r \rfloor = 0$. Telescoping from $m=0$ instead gives
$L_r(0,d) = \sum_m c^{(m)}(d) = d$, while $f = 2^r-1$ kills every term. $\square$

\medskip
\par\noindent
{\bf Lemma 3 (monotonicity in frequency).}
{\em For fixed $d$, the lag $L_r(f,d)$ is non-increasing in $f$. More precisely, if $f<f'$
and $i_0$ is the highest bit in which they differ, then}
\begin{equation}
L_r(f,d) - L_r(f',d) \;\ge\; d_{\,r-1-i_0} \;\ge\; 0.
\label{eq:dyadic_monotone}
\end{equation}

\medskip
\par\noindent
{\em Proof.} Write $m_0 = r-1-i_0$; since $f<f'$ we have $f_{i_0}=0$ and $f'_{i_0}=1$. In
the difference of two copies of (\ref{eq:dyadic_bitwise}), terms with $i>i_0$ cancel, the
$i=i_0$ term contributes $+c^{(m_0)}(d)$, and each term with $i<i_0$ (i.e.\ $m>m_0$)
contributes at least $-c^{(m)}(d)$. Hence
\begin{equation}
L_r(f,d) - L_r(f',d) \;\ge\; c^{(m_0)}(d) - \sum_{m>m_0} c^{(m)}(d) \;=\; d_{m_0}
\end{equation}
by (\ref{eq:dyadic_telescope}). $\square$

\medskip
\par\noindent
Lemma 3 is the algebraic form of an obvious physical statement: tree dedispersion delays the
high tree-freq channels less than the low ones, the total spread being
$L_r(0,d) - L_r(2^r-1,d) = d$. Note that the lag {\em difference} can vanish even when
$f \ne f'$ (when $d_{m_0}=0$ and the lower bits conspire); this is the ``relative lags happen
to cancel'' phenomenon which makes {\tt SparseTile}'s content bits sometimes a proper subset
of the top $(i_0+1)$ bits (\S\ref{sec:avar_dbits}).

\medskip
\par\noindent
{\bf Lemma 4 (the lag function is linear in the DM bits).}
{\em For fixed $f$, the lag function is a non-negative linear form in the binary digits of
$d$:}
\begin{equation}
L_r(f,d) \;=\; \sum_{b=0}^{r-1} d_b \; \lambda_b(f),
  \hspace{1.2cm}
\lambda_b(f) \;\equiv\; L_r\big(f, 2^b\big)
  \;=\; \big(1 - f_{r-1-b}\big) + \sum_{m<b} \big(1-f_{r-1-m}\big) \, 2^{\,b-1-m} \;\ge\; 0.
\label{eq:dyadic_bitlinear}
\end{equation}
{\em Equivalently, $L_r(f,\cdot)$ is additive on bitwise-disjoint DM indices: if
$d \,\&\, e = 0$ then $L_r(f, d\,|\,e) = L_r(f,d) + L_r(f,e)$.}

\medskip
\par\noindent
{\em Proof.} Expanding $\lfloor d/2^{m+1}\rfloor$ in binary,
$c^{(m)}(d) = \lfloor d/2^{m+1}\rfloor + d_m = d_m + \sum_{j>m} 2^{\,j-m-1} d_j$,
which is a linear form in the bits of $d$ with
non-negative integer coefficients; substituting into (\ref{eq:dyadic_bitwise}) and exchanging
the order of summation gives (\ref{eq:dyadic_bitlinear}), and evaluating it at $d = 2^b$
identifies $\lambda_b(f)$. A linear form in the bits is additive whenever the bits add, i.e.\
whenever the two arguments are bitwise disjoint. $\square$

\medskip
\par\noindent
Physically: each DM {\em bit} contributes its own delay ramp $\lambda_b(\cdot)$ across the
band, and a general DM index simply superposes the ramps of its set bits. (The lag is
{\em not} linear in $d$ itself: $\lambda_b(f) \ne 2^b \lambda_0(f)$ in general, the difference
being the rounding built into $\Omega_{-1}=1$.)

\medskip
\par\noindent
{\bf Lemma 5 (bit spreading).}
{\em Write $\Delta_{uv}(d) \equiv L_K(u,d) - L_K(v,d)$. Then for all $u \le v$ and all
bitwise-disjoint $d, e$,}
\begin{equation}
\Delta_{uv}\big( d \,|\, e \big) \;=\; \Delta_{uv}(d) + \Delta_{uv}(e) \;\ge\; \Delta_{uv}(d)
  \;\ge\; 0 .
\label{eq:dyadic_spread}
\end{equation}
{\em Consequently $d \subseteq d'$ implies $\Delta_{uv}(d') \ge \Delta_{uv}(d) \ge 0$: setting
DM bits can only increase every pairwise separation of the impulse response
(\ref{eq:dyadic_impulse}).}

\medskip
\par\noindent
{\em Proof.} The equality is Lemma 4 applied to $u$ and to $v$ and subtracted. Both
inequalities are Lemma 3: $\Delta_{uv}(x) \ge 0$ for $u \le v$ and any $x$, applied to
$x = e$ and to $x = d$. For the consequence, take $e = d' \setminus d$ (bitwise), so that
$d' = d \,|\, e$ with $d \,\&\, e = 0$. $\square$

\medskip
\par\noindent
Three remarks.

First, {\em when is the increment zero?} By (\ref{eq:dyadic_bitwise}),
$\Delta_{uv}(2^b) = \sum_{m \le b} \alpha_m c^{(m)}(2^b)$ with
$\alpha_m = v_{K-1-m}-u_{K-1-m}$, since $c^{(m)}(2^b) = 0$ for $m > b$. If $i_0$ is the
highest bit in which $u$ and $v$ differ and $m_0 = K-1-i_0$, then $\alpha_m = 0$ for $m<m_0$,
so $\Delta_{uv}(2^b) = 0$ whenever $b < m_0$, while $\Delta_{uv}(2^{m_0}) \ge 1$ by
(\ref{eq:dyadic_monotone}). (For $b > m_0$ the increment can also vanish accidentally, by the
cancellation noted after Lemma 3.) So a low DM bit $b$ is felt only by pairs of channels lying
in different coarse-frequency blocks of size $2^{K-b}$; a pair of nearby channels sees no
change at all. This is why, numerically, most instances of (\ref{eq:bitwise_property}) are
exact equalities, increasingly so for the low DM bits: the gridding weights $w_u$ of one input
channel are supported on a short contiguous run of tree-freq channels.

Second, Lemma 5 contains the ``dyadic'' statement that within an aligned dyadic block of DM
indices every pairwise separation is smallest at the first index of the block: that is the
case $(d,e) = (2^kj,\, i)$ with $0 \le i < 2^k$.

Third, the same dyadic statement can be read off the lag-additivity identity of
\ddnotes{The splitting identity} with $r''=k$, $r'=K-k$: writing $u = 2^{K-k}\bar u + u'$,
\begin{equation}
L_K\big(u,\, 2^k j + i\big) \;=\; L_{K-k}(u', j) \;+\; \big(2^k - 1 - \bar u\big) \, j
   \;+\; L_k(\bar u, i),
\label{eq:dyadic_split}
\end{equation}
whose only $i$-dependent term $L_k(\bar u, i)$ depends on $u$ only through its coarse block
index $\bar u$, and is non-increasing in $\bar u$ by Lemma 3. In words: the low DM bits
perform a rank-$k$ dedispersion {\em of the coarse blocks}, sliding them apart while leaving
the arrangement of channels within each block untouched. This is a useful picture, but it is
not needed for the proof: Lemma 4 supersedes it.

\subsection{Spreading out decreases the variance}
\label{ssec:dyadic_kernel}

\par\noindent
{\bf Lemma 6 (kernel autocorrelations are unimodal).}
{\em For every peak-finding profile $p$ (\ddnotes{Peak-finding}), the autocorrelation
$A_p[\delta] = \sum_t h_p[t] h_p[t+\delta]$ is even in $\delta$ and non-increasing in
$|\delta|$.}

\medskip
\par\noindent
{\em Proof.} Evenness is immediate. Write $B_n$ for the boxcar of length $n$ (all entries
1). Each peak-finding kernel is a non-negative combination of boxcars which share a common
center and whose lengths have a common parity: with $w = 2^\lambda$,
\begin{equation}
h_{\lambda,0} = B_w, \hspace{1cm}
h_{\lambda,1} = B_{2w}, \hspace{1cm}
h_{\lambda,2} = \tfrac12 B_{3w} + \tfrac12 B_w, \hspace{1cm}
h_{\lambda,3} = \tfrac12 B_{4w} + \tfrac12 B_{2w},
\end{equation}
where in each case the boxcars are centered on the center of $h_{\lambda,q}$
(e.g.\ $[\tfrac12, 1, \tfrac12] = \tfrac12 [1,1,1] + \tfrac12 [0,1,0]$). Hence $A_p$ is a
non-negative combination of cross-correlations
$\phi_{mn}[\delta] = \sum_t B_m[t] \, B_n[t+\delta]$ of co-centered boxcars, and for
$m \equiv n$ (mod 2) and integer $\delta$,
\begin{equation}
\phi_{mn}[\delta] \;=\; \max\left( 0, \; \min\Big( m, \; n, \; \frac{m+n}{2} - |\delta| \Big) \right),
\end{equation}
which is even and non-increasing in $|\delta|$. A non-negative combination of such functions
has the same property. $\square$

\subsection{The theorem}
\label{ssec:dyadic_theorem}

\par\noindent
{\bf Theorem 7.}
{\em Let $K \ge 0$, let $w_u \ge 0$ for $0 \le u < 2^K$, and let
$\phi : \mathbb{Z} \rightarrow \mathbb{R}$ be even and non-increasing on $\delta \ge 0$.
Define}
\begin{equation}
A[d] \;=\; \sum_{u,v} w_u \, w_v \; \phi\big( L_K(u,d) - L_K(v,d) \big),
  \hspace{1.5cm} 0 \le d < 2^K.
\label{eq:dyadic_abstract}
\end{equation}
{\em Then $A[d'] \le A[d]$ whenever $d \subseteq d'$ bitwise. In particular
$A[2^kj] \ge A[2^kj+i]$ for all $0 \le k \le K$, $0 \le j < 2^{K-k}$, $0 \le i < 2^k$, and
$A[0] \ge A[d]$ for every $d$.}

\medskip
\par\noindent
{\em Proof.} Let $d \subseteq d'$. By Lemma 5, for every pair $u \le v$ we have
$0 \le \Delta_{uv}(d) \le \Delta_{uv}(d')$, whence $\phi(\Delta_{uv}(d)) \ge
\phi(\Delta_{uv}(d'))$ because $\phi$ is even and non-increasing on non-negative arguments.
Multiply by $w_uw_v \ge 0$ and sum over pairs (the sum over $u>v$ duplicates the sum over
$u<v$, since $\phi$ is even and $\Delta_{vu} = -\Delta_{uv}$). The dyadic-block statement is
the case $d = 2^kj$, $d' = 2^kj+i$, whose bits are disjoint. $\square$

\medskip
\par\noindent
{\bf Corollary 8.}
{\em The bitwise monotonicity (\ref{eq:bitwise_property}), and hence the dyadic-block property
(\ref{eq:dyadic_property}), holds for every column of the variance map
$A[d_m, \mbox{\rm subband}, p, F]$ of \S\ref{ssec:svd_exact_map}, in the no-detrender setting
of \S\ref{ssec:dyadic_setup}, for every {\tt DedispersionTree} (any $\gamma$, any $e$), every
subband (any pf\_level $l$, both the aligned Case 1 and the half-aligned Case 2), every
profile $p$, and every input channel $F$.}

\medskip
\par\noindent
{\em Proof.} Lemma 1 writes the column in the form (\ref{eq:dyadic_abstract}) with
$w_u \ge 0$; Lemma 6 gives $\phi = A_p$ the required monotonicity; apply Theorem 7. For
$\gamma>0$ the exported array is $A[d_m] = \tilde A[d_m + 2^{K-1}]$, where $\tilde A$
satisfies the theorem; since $d_m \subseteq d'_m$ (with both $< 2^{K-1}$) implies
$d_m + 2^{K-1} \subseteq d'_m + 2^{K-1}$, the property is inherited. $\square$

\medskip
\par\noindent
Three extensions are immediate, because bitwise monotonicity is preserved by non-negative
combinations and by pointwise maxima over a common DM index:
\begin{itemize}
\item {\bf Summing over frequency channels.} $y[d_m] = \sum_F V_F \, A[d_m,\cdot,\cdot,F]$
has the property for any $V_F \ge 0$; this is the array {\tt PfAvarExact.tree\_variance}.
\item {\bf Maximizing over profiles or multiplets.} Arrays which share the same $d_m$ index
(e.g.\ the $P$ profiles) inherit the property under sums and pointwise maxima, since the
maximizing index is common to all of them. The multiplets of one subband instead share the
{\em coarse} DM index $d_{\rm hi}$, and the property passes to them too, because $d_{\rm hi}$
and $d_{\rm lo}$ occupy disjoint bit ranges of $d_m = 2^l d_{\rm hi} + d_{\rm lo}$: each
individual multiplet's array $d_{\rm hi} \mapsto A[2^l d_{\rm hi} + d_{\rm lo}]$ is bitwise
monotone, since $d_{\rm hi} \subseteq d'_{\rm hi}$ implies
$2^l d_{\rm hi} + d_{\rm lo} \subseteq 2^l d'_{\rm hi} + d_{\rm lo}$; and the array the
peak-finder actually forms,
$B[d_{\rm hi}] = \max_{d_{\rm lo}} A[2^l d_{\rm hi} + d_{\rm lo}]$, equals
$A[2^l d_{\rm hi}]$ by Theorem 7 (as $0 \subseteq d_{\rm lo}$), i.e.\ it is the fine array
sampled at multiples of $2^l$, and is therefore bitwise monotone in $d_{\rm hi}$. (Different
subbands have different $l$ and hence different DM ranges, so a maximum over subbands only
makes sense once they are put on a common DM grid; on such a grid it inherits the property
as well.)
\item {\bf Coarse-graining the DM axis.} If the DM axis is coarse-grained by a factor $2^c$
using max (as in the dedispersion output arrays, \ddnotes{Dedispersion output arrays}),
the coarse-grained array equals the fine array sampled at $d_m \equiv 0 \pmod {2^c}$, and
$i \subseteq i'$ implies $2^c i \subseteq 2^c i'$, so the coarse array inherits the property.
\end{itemize}
The argument also covers the {\em approximate} variance arrays of {\tt PfAvarApproximation}
(\S\ref{sec:variance_svd}): there each per-band contribution is again of the form
(\ref{eq:dyadic_closedform}) (with $K$ the rank of the corresponding sub-block tile), and the
reported variance is a non-negative average of such contributions over bands, all sharing one
DM index.

\subsection{Sharpness}
\label{ssec:dyadic_sharpness}

Neither hypothesis of Theorem 7 can be dropped.
\begin{itemize}
\item {\em Non-negative weights.} Take $K=1$, $h=[1]$ (profile $p=0$) and $w = (1,-1)$. Then
$A[0]=0$ and $A[1]=2$, so even the two-element instance $0 \subseteq 1$ of
(\ref{eq:bitwise_property}) fails. This is why the argument does not survive a
{\tt Detrender2d} in front of the dedisperser: the effective per-channel response is no longer
a non-negative weight times a single delta, so neither hypothesis of Theorem 7 survives. This
is not a gap in the proof but a real failure: computing the detrended variance map by brute
force (\S\ref{ssec:variance_map_bruteforce}) at rank 6 with {\tt subband\_counts}=[2,2,1],
about half of all aligned dyadic blocks of $A[\cdot,\mbox{subband},p,F]$ are violated, with
block maxima up to 5 times the block-bottom value for an $(n,W)=(2,4)$ detrender. For the
$W=0$ detrender the failure is extreme and systematic: the detrender annihilates the zero-DM
mode, so $A[0]$ is close to zero and is the block {\em minimum} rather than its maximum.
\item {\em Unimodal kernel autocorrelation.} Take $K=3$, $w=(1,1,1,1,1,1,1,1)$ (a full-band,
uniform gridding footprint) and the two-hump kernel $h = [1,0,0,0,1]$, whose autocorrelation
$A_h = [2,0,0,0,1]$ is not monotone. Then $A = (128,64,48,32,30,32,26,24)$, which violates
(\ref{eq:bitwise_property}) at $4 \subseteq 5$ ($A[5] = 32 > A[4] = 30$), and violates the
dyadic-block property on the block $\{4,5\}$. All four current profile shapes are co-centered
sums of boxcars, so Lemma 6 applies; but a future profile which is not (for example a matched
filter for a scattered or multi-component pulse) would invalidate the conclusion, not merely
the proof -- and, as this example shows, for perfectly ordinary gridding weights.
\end{itemize}

Nor can the bitwise-subset order be enlarged: it is exactly the set of pairs which the
hypotheses of Theorem 7 constrain.

\medskip
\par\noindent
{\bf Proposition 9 (the order is exactly right).}
{\em Fix $K$ and $0 \le d, d' < 2^K$. The following are equivalent:}
\begin{itemize}
\item[(i)] $A[d'] \le A[d]$ for every $w \ge 0$ and every even, non-increasing $\phi$ in
(\ref{eq:dyadic_abstract});
\item[(ii)] $\Delta_{uv}(d') \ge \Delta_{uv}(d)$ for every $u \le v$;
\item[(iii)] $d \subseteq d'$.
\end{itemize}

\medskip
\par\noindent
{\em Proof.} (iii)$\Rightarrow$(ii) is Lemma 5, and (ii)$\Rightarrow$(i) is the proof of
Theorem 7.

(i)$\Rightarrow$(ii): the case $u=v$ is trivial, so let $u<v$, and take the two-channel weight
vector with $w_u = w_v = 1$ and all other entries zero, so that
$A[x] = 2\phi(0) + 2\phi(\Delta_{uv}(x))$; take $\phi(\delta) = 1$ for $|\delta| \le t$ and
$\phi(\delta) = 0$ otherwise, which is even and non-increasing on $\delta \ge 0$. If
$\Delta_{uv}(d') < \Delta_{uv}(d)$, then $t = \Delta_{uv}(d')$ gives $A[d'] > A[d]$.

(ii)$\Rightarrow$(iii): write $\epsilon_m = c^{(m)}(d') - c^{(m)}(d)$ and
$D_m = \lfloor d'/2^m \rfloor - \lfloor d/2^m \rfloor$, so that $\epsilon_m = D_m - D_{m+1}$
and $D_K = 0$ by (\ref{eq:dyadic_telescope}). By (\ref{eq:dyadic_bitwise}),
$\Delta_{uv}(x) = \sum_m \alpha_m c^{(m)}(x)$ with $\alpha_m = v_{K-1-m} - u_{K-1-m}$; for a
pair $u \le v$ whose highest differing bit is $i_0 = K-1-m_0$ we have $\alpha_m = 0$ for
$m<m_0$, $\alpha_{m_0}=1$, and $\alpha_m \in \{-1,0,1\}$ arbitrary and independent for
$m>m_0$ (the lower bits of $u$ and $v$ are unconstrained). Hence (ii) says exactly
\begin{equation}
\epsilon_{m_0} \;\ge\; \sum_{m>m_0} |\epsilon_m|
  \hspace{1.5cm} \mbox{for every } 0 \le m_0 < K .
\end{equation}
Since the right side is $\ge 0$, every $\epsilon_m \ge 0$, so the constraint becomes
$\epsilon_{m_0} \ge \sum_{m>m_0}\epsilon_m = D_{m_0+1}$, i.e.\ $D_{m_0} \ge 2 D_{m_0+1}$. But
$D_{m} - 2D_{m+1} = d'_{m} - d_{m}$, so this says $d'_m \ge d_m$ for every $m$, which is
$d \subseteq d'$. $\square$

\medskip
\par\noindent
In particular $A$ need not be monotone in $d_m$ as an integer, and no weaker hypothesis on the
DM indices will do. A minimal example within the hypotheses of Theorem 7, using a real
peak-finding kernel: $K=2$, $w = (0,1,1,0)$, $h = B_2 = [1,1]$ (so $\phi = A_h = [2,1,0,\dots]$)
gives $A = (8,6,8,6)$, which increases from $d=1$ to $d=2$. The reason is visible in
(\ref{eq:dyadic_bitwise}): channels $u=1$ and $v=2$ have lags $(1,0)$ at $d=1$ but $(1,1)$ at
$d=2$, so the pair is {\em less} spread out at the larger DM index -- while at $d=3 = 1|2$,
which is above both in the subset order, the separation is back to 1. The same effect is easy
to see in the real variance map: for instance in {\tt configs/dedispersion/toy.yml}, tree 0,
subband $n=0$ ($l=1$, $\rho=9$), channel $F=328$, profile $p=0$, one finds
$A[255] = 5.08\times10^{-4} < A[256] = 6.71\times10^{-4}$, with
$A[255 \,\&\, 256] = A[0] = 9.77\times10^{-4}$ above both.

Finally, the statement is about the {\em natural} (non-bit-reversed) DM index
$d_m = 2^l d_{\rm hi} + d_{\rm lo}$, which is the convention used by {\tt PfAvarExact} and
{\tt VarianceMapExact}. Under a bit-reversed DM index, aligned dyadic blocks become strided
sets, and (\ref{eq:bitwise_property}) becomes bitwise monotonicity with respect to the
reversed bit labelling (the subset order is preserved by any relabelling of the bits, so the
theorem is convention-independent; only its ``dyadic block'' phrasing is not).

\clearpage

\section{{\tt PfAvarExact} and {\tt VarianceMapExact}}
\label{sec:analytic_variance}

The second step of peak-finding (\ddnotes{Peak-finding}) normalizes its output to ``sigmas'',
by applying weights which are the inverse square roots of the output variances.
This appendix describes the two classes which compute them, both in
\verb|pirate_frb/slow_avar/|:
\begin{itemize}
\item \verb|PfAvarExact| (\S\ref{sec:avar_statement}--\S\ref{sec:avar_validation}) computes the
variance of the dedispersion outputs, given the per-channel variances of the dedisperser input
array. In the language of \S\ref{sec:variance_map}, Eq.\ (\ref{eq:variance_map_def}), it
evaluates the matrix-vector product $y = Av$ for one specific $v$, without ever forming $A$.
\item \verb|VarianceMapExact| (\S\ref{ssec:svd_exact_map}--\S\ref{ssec:svd_exact_cost}) takes a
\verb|PfAvarExact| and compresses the map $A$ itself into a canonical low-rank (SVD) form. This
is worth doing because in operation $v$ drifts (gain drift, RFI excision), and we want to
recompute the peak-finding variances cheaply whenever it is updated.
\end{itemize}
Both use the helper classes \verb|SparseTile|, \verb|SparseTileTriple|, \verb|PfVariance| and
\verb|PfVarianceConvolver|.\footnote{This is an exact but slow reference
implementation. Production uses a faster approximate algorithm ({\tt PfAvarApproximation},
which is also ported to C++); see \S\ref{sec:variance_svd}.}

One difference from \S\ref{ssec:variance_map_setup} applies throughout: there, the underlying
linear operator includes a {\tt Detrender2d} before tree dedispersion, whereas here we consider
``bare'' tree dedispersion. So the matrix called $A$ in this appendix is the $A$ of
Eq.\ (\ref{eq:variance_map_def}) with the detrending step omitted; extending the construction to
include the detrender is future work.

To keep the presentation manageable, \S\ref{sec:avar_statement}--\S\ref{sec:avar_validation}
make two simplifying assumptions: we consider
only the ``base tree'' (no time downsampling and no early triggers, \ddnotes{Time-downsampled trees and early triggers}),
and we disable subbands (i.e.\ $R=0$ and $C_l = \{1\}$, \ddnotes{Subband search}).
Then there is a single multiplet ($M=1$), and the output of the calculation is a shape $(2^r,P)$
array, indexed by delay $0 \le d < 2^r$ and peak-finding profile $0 \le p < P$, where
$r = {\tt toplevel\_tree\_rank}$. These assumptions are expositional only -- the code handles the
general case -- and \S\ref{ssec:svd_exact_map} onwards drops them.

\subsection{Statement of the problem}
\label{sec:avar_statement}

The dedisperser input is a shape $(N_{\rm freq}, N_{\rm time})$ array ${\tt in}[F,t]$.
We assume that its elements are uncorrelated, with a variance which depends on frequency but not
time:
\begin{equation}
\big\langle {\tt in}[F,t] \, {\tt in}[F',t'] \big\rangle = V_F \, \delta_{FF'} \, \delta_{tt'}
  \label{eq:avar_input}
\end{equation}
(In the code, the length-$N_{\rm freq}$ array $V_F$ is the constructor argument
\verb|freq_variances|.)
The three stages downstream -- tree gridding (\ddnotes{Tree gridding kernel}), tree dedispersion
(\ddnotes{Tree dedispersion}), and convolution with the peak-finding kernels $h_p$
(\ddnotes{Peak-finding}) -- are all linear, and are all time-translation invariant, so their
composition can be written as
\begin{equation}
y_p[d,t] = \sum_{F=0}^{N_{\rm freq}-1} \sum_{t'} \big( h_p * G_{Fd} \big)[t-t'] \; {\tt in}[F,t']
\end{equation}
Here, $G_{Fd}$ is the {\em impulse response}: the tree dedispersion output at delay $d$, if
gridding and $\DD(r)$ are run on an input array whose only nonzero element is ${\tt in}[F,0]=1$.
Combining with Eq.\ (\ref{eq:avar_input}), the array we want to compute is
\begin{equation}
{\rm Var}[d,p] \equiv \big\langle y_p[d,t]^2 \big\rangle
  = \sum_F V_F \, \big\| h_p * G_{Fd} \big\|^2
  \hspace{1.5cm} \mbox{where } \|g\|^2 \equiv \sum_t g[t]^2
  \label{eq:avar_master}
\end{equation}
Note that the RHS does not depend on $t$, i.e.\ every output time sample has the same
variance.\footnote{This is a steady-state statement. Near the beginning of a run, output samples
have a {\em smaller} variance, since part of the impulse response reaches back to times before the
start of the stream (recall from \ddnotes{Tree dedispersion} that ${\tt in}[\cdots,t]$ is
defined to be zero for $t<0$). The {\tt DedispersionPlan} knows when each DM row reaches steady
state ({\tt compute\_steady\_state\_it0}), and the Monte Carlo check of \S\ref{sec:avar_validation}
uses this.}

We can make Eq.\ (\ref{eq:avar_master}) completely explicit. Recall from
\ddnotes{Non-recursive formulation of tree dedispersion} that $\DD(r)$ maps a one-hot
tree-freq channel $f$ to a single
sample, at lag $L_r(f,d)$ with amplitude $2^{-r/2}$, and from \ddnotes{Tree gridding kernel} that input
channel $F$ is gridded onto tree-freq channels with weights $w_{fF}$. Therefore
\begin{equation}
G_{Fd}[t] = 2^{-r/2} \sum_{f \in {\mathcal S}_F} w_{fF} \; \delta_{t, L_r(f,d)}
  \label{eq:avar_impulse}
\end{equation}
where the {\em gridding footprint} ${\mathcal S}_F \equiv \{ f : w_{fF} > 0 \}$ is the set of
tree-freq channels whose interval $I_f$ overlaps the input channel $[F,F+1)$. Since the intervals
$I_f$ tile the band, ${\mathcal S}_F$ is a contiguous range
\begin{equation}
f_0(F) \le f < f_1(F)
\end{equation}
which is usually just a few channels wide (in the CHORD config, $N_{\rm freq} = 28160$ input
channels are gridded onto $2^r = 65536$ tree-freq channels).

Thus the impulse response of a single input channel is a handful of delta functions, and
Eqs.\ (\ref{eq:avar_master}), (\ref{eq:avar_impulse}) are a complete specification of the answer.
However, evaluating them directly, for each $(F,d)$ pair, would be too slow
($N_{\rm freq} 2^r \approx 2 \times 10^9$ pairs in the CHORD config), and would also be wasteful:
we will see that the $2^r$ delays give rise to far fewer than $2^r$ distinct answers.
The rest of this section describes the algorithm which is actually used.

\subsection{Variance from autocorrelations}
\label{sec:avar_convolver}

The lowest-level building block is the map $g \mapsto \|h_p * g\|^2$, for a short time series $g$
and each of the $P$ peak-finding kernels. Expanding the convolutions and collecting terms by lag,
\begin{equation}
\big\| h_p * g \big\|^2 = \sum_{\delta} \rho_g[\delta] \, A_p[\delta]
  = \rho_g[0] \, A_p[0] + 2 \sum_{\delta \ge 1} \rho_g[\delta] \, A_p[\delta]
  \label{eq:avar_autocorr}
\end{equation}
where $\rho_g$ and $A_p$ are the (one-sided) autocorrelations of the data and of the kernel:
\begin{equation}
\rho_g[\delta] = \sum_t g[t] \, g[t+\delta]
  \hspace{1.5cm}
A_p[\delta] = \sum_t h_p[t] \, h_p[t+\delta]
\end{equation}
Two features of Eq.\ (\ref{eq:avar_autocorr}) are worth noting.
\begin{itemize}
\item $A_p[\delta]$ vanishes for $\delta \ge {\rm len}(h_p)$, so lags
$\delta \ge T_{\rm max} \equiv \max_p {\rm len}(h_p) = 2 W_{\rm max}$ never contribute.
The class \verb|PfVarianceConvolver| precomputes the shape $(P,T_{\rm max})$ table $A_p[\delta]$
once, and Eq.\ (\ref{eq:avar_autocorr}) then becomes a small matrix multiply which returns all $P$
profiles at once. The $\delta=0$ column is just the kernel norms:
\begin{equation}
A_p[0] = \|h_p\|^2 = 2^\lambda \times \big\{ 1, \; 2, \; \tfrac{3}{2}, \; \tfrac{5}{2} \big\}_q
\end{equation}
where $p = 3\lambda + q$ as in \ddnotes{Peak-finding}.
Evaluating (\ref{eq:avar_autocorr}) for a time series whose support has length $n_t$ costs
$O(n_t \min(n_t, T_{\rm max}))$ for $\rho_g$, plus $O(\min(n_t,T_{\rm max}) \, P)$ for the
contraction.

\item The LHS is invariant under a time translation $g[t] \rightarrow g[t-T]$. Therefore, we only
need each impulse response {\em up to an arbitrary time shift}. This is the key to the
compression in \S\ref{sec:avar_dbits}--\ref{sec:avar_sparsetile}: most of the $d$-dependence of
$G_{Fd}$ is a pure time shift, and can simply be discarded.
\end{itemize}

\subsection{Which delay bits matter}
\label{sec:avar_dbits}

By the time-translation invariance just noted, only the {\em relative} lags within the footprint
matter in Eq.\ (\ref{eq:avar_impulse}). From the definition of the lag function
(\ddnotes{Non-recursive formulation of tree dedispersion}),
the relative lag of two tree-freq channels $f,f'$ is
\begin{equation}
L_r(f,d) - L_r(f',d) = \sum_{i=0}^{r-1} \sum_{j=0}^{r-1} \Omega_{i+j-r} \, (f'_i - f_i) \, d_j
  \label{eq:avar_lagdiff}
\end{equation}
Suppose that $f$ and $f'$ agree in all bits above position $i_0$, so that only terms with
$i \le i_0$ survive in Eq.\ (\ref{eq:avar_lagdiff}). Since $\Omega_n = 0$ for $n < -1$, a term also
vanishes unless $j \ge r-1-i$. Therefore the relative lag depends only on the bits $d_j$ with
$j \ge r-1-i_0$, i.e.\ on the top $(i_0+1)$ bits of $d$.
For a contiguous footprint, all channels agree in the bits above
\begin{equation}
i_0(F) \equiv \Big\lfloor \log_2 \big( f_0(F) \oplus (f_1(F)-1) \big) \Big\rfloor
  \label{eq:avar_i0}
\end{equation}
where $\oplus$ denotes bitwise XOR (and we define $i_0=-1$ if the footprint is a single channel).
We conclude:
\begin{quote}
The contribution of input channel $F$ to Eq.\ (\ref{eq:avar_master}) depends on the delay $d$ only
through its top $(i_0(F)+1)$ bits.
\end{quote}
In particular, if the footprint is a single tree-freq channel $f$, then the contribution is
$d$-independent, and Eqs.\ (\ref{eq:avar_master}), (\ref{eq:avar_impulse}) collapse to
\begin{equation}
V_F \big\| h_p * G_{Fd} \big\|^2 = 2^{-r} \, V_F \, w_{fF}^2 \, \|h_p\|^2
\end{equation}

Eq.\ (\ref{eq:avar_i0}) has a simple interpretation: $i_0(F)$ is the largest $i$ such that the
footprint straddles a boundary between two ``coarse'' frequency channels of width $2^i$
(cf.\ the coarse-freq blocking of \ddnotes{Tree dedispersion}). Equivalently, $\DDS(i_0)$ is
the last dedispersion step which
combines two nonzero halves of the footprint with each other. There are only $2^{r-i}$ such
boundaries in the band, and each one is straddled by at most one input channel, so at most
$O(2^{r-i})$ input channels have $i_0 \ge i$: large $i_0(F)$ is rare, but (as we will see)
expensive.

The algorithm below never computes $i_0(F)$ explicitly. Instead, it discovers the same structure
-- and sometimes a strictly smaller set of relevant bits, when relative lags happen to cancel --
while propagating the impulse response through the dedispersion steps.

\subsection{Sparse propagation of the impulse response}
\label{sec:avar_sparsetile}

Recall from \ddnotes{Tree dedispersion} that the intermediate array after $k$ dedispersion steps
has shape $(2^{r-k}, 2^k, N_{\rm time})$, with axes (coarse-freq, delay, time). For a one-hot
input, such an array is highly structured, and the class \verb|SparseTile| represents it with the
following members:
\begin{itemize}
\item $f_0, n_f$: the array vanishes outside coarse-freq channels $f_0 \le f < f_0 + n_f$.
\item $n_t$: the stored data vanishes outside time samples $0 \le t < n_t$.
\item {\tt dbits}: a bitmask selecting a subset of the $k$ delay bits (the ``content bits'').
We write $|{\tt dbits}|$ for its popcount.
\item {\tt data}: an array of shape $(n_f, 2^{|{\tt dbits}|}, n_t)$.
\item $t_0$ and $\tau_0, \cdots, \tau_{k-1}$: a constant time shift, and one time shift per delay bit.
\item $s$: a scalar multiplier.
\end{itemize}
The dense array represented by the tile is
\begin{equation}
{\tt arr}[f,d,t] = s \cdot {\tt data}\big[ f - f_0, \; d|_{\tt dbits}, \; t - T(d) \big]
  \hspace{1.2cm}
T(d) = t_0 + \sum_{i=0}^{k-1} d_i \, \tau_i
  \label{eq:avar_tile}
\end{equation}
where $d|_{\tt dbits}$ denotes the selected bits of $d$, packed (highest bit first) into an index
$0 \le d|_{\tt dbits} < 2^{|{\tt dbits}|}$.
The compression is therefore two-fold: the array is stored only on its $(n_f, n_t)$ support, and
its $2^k$ delay rows are represented by $2^{|{\tt dbits}|}$ stored rows, plus a time shift $T(d)$
which is affine in the bits of $d$.

The starting point ($k=0$) is the gridding output for a one-hot input, which is a tile with
$f_0 = f_0(F)$, $n_f = f_1(F) - f_0(F)$, $n_t = 1$, ${\tt dbits} = 0$, $s=1$, no shifts, and
${\tt data}[f,0,0] = w_{fF}$.

The tile representation is preserved by a dedispersion step $\DDS(k)$. Consider an output
coarse-freq channel $f$, whose inputs are the ``lower'' channel $2f$ and the ``upper'' channel
$2f+1$, and write the output delay index as $d' = 2d+b$ (so that $d'_0 = b$ and $d'_{j+1} = d_j$).
There are three cases:
\begin{itemize}

\item {\bf Upper half only} (the lower half is outside the footprint).
Then ${\tt out}[f,2d+b,t] = 2^{-1/2} \, {\tt in}[2f+1,d,t]$ for both $b=0,1$, i.e.\ the output does
not depend on the new delay bit at all. The tile is unchanged, except that its bitmask and shift
vector are ``lifted'' by one bit,
\begin{equation}
{\tt dbits} \rightarrow 2 \, {\tt dbits}
  \hspace{1.5cm}
(\tau_0, \cdots, \tau_{k-1}) \rightarrow (0, \tau_0, \cdots, \tau_{k-1})
  \label{eq:avar_lift}
\end{equation}
and $s \rightarrow s/\sqrt{2}$. No data is touched.

\item {\bf Lower half only}. Then ${\tt out}[f,2d+b,t] = 2^{-1/2} \, {\tt in}[2f,d,t-d-b]$, i.e.\ the
output depends on the new delay bit only through a time shift. The bitmask and shift vector are
lifted as in Eq.\ (\ref{eq:avar_lift}), and in addition the dedispersion lag $(d+b)$ is absorbed
into the shifts, by adding
\begin{equation}
(1, 1, 2, 4, \cdots, 2^{k-1})
  \label{eq:avar_ddshifts}
\end{equation}
to the (lifted) length-$(k+1)$ shift vector, using $d + b = b + \sum_{j \ge 1} d'_j 2^{j-1}$.
Again no data is touched.

\item {\bf Both halves.} Now the output is a sum of two ``legs'', each of which is a lifted copy of
one input tile, and each of which carries its own shift vector: $\sigma^U$ for the upper leg, and
$\sigma^L$ for the lower leg (which includes the extra dedispersion lags of
Eq.\ \ref{eq:avar_ddshifts}). Only the {\em difference} between the two shift vectors is physical,
so we factor the common part $\tau_j = \min(\sigma^L_j, \sigma^U_j)$ out into the output tile's
shifts, and fold the residuals $(\sigma^L_j - \tau_j)$ and $(\sigma^U_j - \tau_j)$ into the data,
by materializing each output row and adding the two legs at their residual offsets (the constant
shifts $t_0$ are treated the same way). A delay bit $j$ becomes a content bit of the output if the
two legs' residuals differ ($\sigma^L_j \ne \sigma^U_j$), or if it was already a content bit of
either leg.

This is the only case in which content bits are created, data is copied, and the time support
grows (by up to $\sum_j (\sigma^L_j - \tau_j) \le 2^k$).

\end{itemize}
Since the footprint is contiguous, at most two coarse-freq channels -- the first and the last --
can have a missing half. The class \verb|SparseTileTriple| exploits this by splitting the array
into (at most) three tiles: first channel, bulk, last channel. The two edge tiles can then take
the cheap single-leg cases, while the bulk is processed with the both-halves case uniformly.

After $r$ dedispersion steps, only one coarse-freq channel is left, and the resulting tile
($k=r$, $n_f=1$) is an exact, compressed representation of the impulse response $G_{Fd}$ of
Eq.\ (\ref{eq:avar_impulse}). Its content bits are always a subset of the top $(i_0(F)+1)$ bits
predicted in \S\ref{sec:avar_dbits}, and are sometimes a proper subset, since the merge rule above
only marks a bit when the two legs' residual shifts actually differ.

\subsection{Accumulating over frequency}
\label{sec:avar_pfvariance}

Let ${\tt data}$, {\tt dbits} and $s$ be the members of the final ($k=r$) tile for input channel
$F$. Since $\|h_p * g\|^2$ is invariant under time translation, the shifts $T(d)$ in
Eq.\ (\ref{eq:avar_tile}) drop out of the variance entirely, leaving
\begin{equation}
\big\| h_p * G_{Fd} \big\|^2 = s^2 \, \big\| h_p * {\tt data}[0, \, d|_{\tt dbits}, \, \cdot] \big\|^2
  \label{eq:avar_tile_variance}
\end{equation}
So we simply run each of the $2^{|{\tt dbits}|}$ stored rows through Eq.\ (\ref{eq:avar_autocorr}).
The result is a shape $(2^{|{\tt dbits}|}, P)$ array containing the entire contribution of input
channel $F$ to Eq.\ (\ref{eq:avar_master}).

The frequency sum in Eq.\ (\ref{eq:avar_master}) is then accumulated in factored form. A
\verb|PfVariance| object is a small dictionary which maps a bitmask to an array of shape
$(2^{|{\tt dbits}|}, P)$, and represents
\begin{equation}
{\rm Var}[d,p] = \sum_{\rm terms} {\tt arr}^{({\tt dbits})}\big[ d|_{\tt dbits}, \, p \big]
  \label{eq:avar_terms}
\end{equation}
Accumulating channel $F$ means scaling its array by $V_F$, and adding it to the term with the same
bitmask (creating a new term if there is none). Only at the very end is
Eq.\ (\ref{eq:avar_terms}) ``unpacked'' into the dense shape $(2^r,P)$ output array. Note that the
largest object ever materialized is a single channel's tile: the full
$(N_{\rm freq}, 2^r, N_{\rm time})$ impulse response never exists in memory.

The cost is dominated by input channels with large $i_0(F)$, whose final tile has $O(2^{i_0})$
stored rows, each of length $n_t = O(2^{i_0})$. As noted in \S\ref{sec:avar_dbits}, such channels
are rare -- at most $O(2^{r-i})$ input channels have $i_0 \ge i$ -- but the few channels which
straddle a high-level boundary (e.g.\ the middle of the band) are expensive enough that this exact
calculation is only practical at moderate ranks. This is the main reason that it is used as a
testing/reference implementation, rather than in production.

\subsection{Validating the analytic variances}
\label{sec:avar_validation}

The pieces above are unit-tested against dense reference implementations
(\verb|pirate_frb test --avar|): the \verb|SparseTile| operations are compared to a dense
$\DDS(k)$ applied to randomly generated tiles, tree gridding is compared to
\verb|ReferenceTreeGriddingKernel|, dedispersion of a one-hot input is compared to
\verb|ReferenceTree|, and \verb|PfVarianceConvolver| is compared both to a brute-force evaluation
of $\|h_p * x\|^2$, and to the kernels which \verb|ReferencePeakFindingKernel| actually uses.

End-to-end, the analytic variances are checked by Monte Carlo
(\verb|pirate_frb check_avar_mc|): Gaussian noise with per-channel variance $V_F$ is fed through a
\verb|ReferenceDedisperser| (with \verb|enable_variances=True| and peak-finding weights set to 1),
and its empirical per-chunk variances are compared to the analytic prediction, restricted to
output channels which have reached steady state.

\clearpage

The rest of this appendix describes \verb|VarianceMapExact|, which compresses the variance map
$A$ itself, instead of evaluating $y = Av$ for one input vector. From here on we drop the two
simplifying assumptions made above: subbands are enabled, and the variance is resolved per
multiplet, exactly as {\tt PfAvarExact} computes it.

\S\ref{sec:variance_svd} is the counterpart of what follows for the {\em approximate} variance
map ({\tt PfAvarApproximation}). The linear algebra is the same in both cases; the differences
are which variance array is compressed, and how finely it resolves the DM axis. There, the
frequency-channel overlaps between the $2^R$ level-0 bands are neglected, so the variance is
resolved per level-0 band $j$ and on only the $(r-L)$ highest DM bits; here it is resolved per
multiplet $m$ and on the $(r-R)$ highest DM bits, with no approximation. One limitation does
carry over from \S\ref{sec:avar_pfvariance}: {\tt PfAvarExact} is practical only at moderate
rank, so {\tt VarianceMapExact} inherits that and cannot be run at CHIME/CHORD scale.

\subsection{The variance map}
\label{ssec:svd_exact_map}

Fix a {\tt DedispersionTree} with rank $r$, peak-finding rank $R$, and $P$ peak-finding
profiles, and write
\begin{equation}
\rho = r-R, \hspace{1.5cm} D = 2^\rho
\end{equation}
for the number of coarse-DM bits and coarse DMs.\footnote{Code names: $r = {\tt tree\_r}$,
$R = {\tt tree\_R}$, $P = {\tt tree\_P}$, $M = {\tt tree\_M}$, all members of
{\tt VarianceMapExact}.}
Recall from \ddnotes{Subband search} that a {\em multiplet} $0 \le m < M$ is a
(frequency subband, fine DM) pair. For each $m$, {\tt PfAvarExact} produces a variance array
\begin{equation}
y_m[d,p], \hspace{1.5cm} 0 \le d < D, \quad 0 \le p < P,
\end{equation}
(the array {\tt tree\_variance[itree][m]}), which we flatten to a length-$DP$ column vector
using the index $\alpha = dP + p$ of \S\ref{ssec:variance_map_setup}. Then, for each
$(\mbox{tree}, m)$, the map $v \mapsto y_m$ is linear:
\begin{equation}
y_m = A^{(m)} v, \hspace{1.5cm} A^{(m)}_{\alpha F} = \mbox{(contribution of channel $F$ to $y_m[d,p]$
per unit input variance)},
\label{eq:svdx_map}
\end{equation}
which is the matrix of Eq.\ (\ref{eq:variance_map_def}), restricted to one multiplet: rows index
outputs, columns index input channels.

Only a contiguous range of input channels contributes,
\begin{equation}
F_{\rm lo} \le F < F_{\rm lo} + n_m,
\label{eq:svdx_support}
\end{equation}
namely those whose gridding footprint ${\mathcal S}_F$ (\S\ref{sec:avar_statement}) intersects
the coarse-freq range of $m$'s frequency subband.\footnote{Code names: $F_{\rm lo} =
{\tt per\_tm\_ifreq\_lo}$, $n_m = {\tt per\_tm\_nf}$, both members of {\tt PfAvarExact}.}
Contiguity holds because the channel map is strictly decreasing, so the ${\mathcal S}_F$ are
intervals which advance monotonically with $F$, and the set of $F$ whose interval meets a fixed
interval is itself an interval. ({\tt PfAvarExact} asserts this at construction.) We restrict
$A^{(m)}$ to those columns from here on, so it is a $(DP) \times n_m$ matrix; it is never
materialized.

Channels at the two ends of the range (\ref{eq:svdx_support}) may overlap the subband only
slightly, and would contribute anomalously small columns to $A^{(m)}$; a fixed relative SVD
threshold would silently truncate them away. We precondition this away with a diagonal matrix
$\Lambda = {\rm diag}(\lambda)$, where $\lambda_F > 0$ is the mean variance contributed by
channel $F$ (defined precisely in the next subsection), and write
\begin{equation}
\bar A = A^{(m)} \Lambda^{-1},
\hspace{1.5cm}
y_m = \bar A \, \Lambda \, v,
\label{eq:svdx_abar}
\end{equation}
so that $\bar A$ has columns of comparable scale. Diagonal scaling is invertible, so this does
not change matrix ranks; its only purpose is to make a fixed relative threshold mean the same
thing for every channel.

\subsection{Factored structure of single-channel variances}
\label{ssec:svd_exact_factored}

{\tt PfAvarExact} already holds the variance contribution of a single channel in a factored
form: {\tt per\_tfm[itree][F][m]} is a {\tt PfVariance} of rank $\rho$
(\S\ref{sec:avar_pfvariance}), i.e.\ a small set of {\em terms}, each keyed by a bitmask $\delta$
over the $\rho$ DM bits (``{\tt dbits}'') and holding an array of shape
$(2^{{\rm popc}(\delta)}, P)$, with
\begin{equation}
y^{(F)}_m[d, p] = \sum_{(\delta, a) \,\in\, {\rm terms}(F,m)} a\big[ {\rm sel}(d, \delta),\, p \big],
\hspace{1.5cm}
y_m = \sum_F v_F \, y^{(F)}_m,
\label{eq:svdx_terms}
\end{equation}
where ${\rm sel}(d, \delta)$ packs the bits of $d$ selected by $\delta$ into an integer
$0 \le {\rm sel}(d, \delta) < 2^{{\rm popc}(\delta)}$ (the function
\verb|SparseTile._remap_d|). That is, a term depends on the DM index $d$ only through the few
bits in $\delta$. Note that the entries of $a$ are variances $\|h_p * g\|^2$, hence
non-negative.

Compared with \S\ref{sec:variance_svd}, the masks $\delta$ here are drawn from the $\rho = r-R$
resolved DM bits rather than $r-L$, and satisfy
${\rm popc}(\delta) \le \min(i_0(F)+1, \rho)$ with $i_0(F)$ as in \S\ref{sec:avar_dbits}. More
channels therefore carry multi-bit masks, and more distinct masks occur per $(\mbox{tree},m)$
pair -- 12 to 25 in {\tt toy.yml}, where the approximation sees only a handful.

The preconditioner in Eq.\ (\ref{eq:svdx_abar}) is defined in terms of the factored form:
\begin{equation}
\lambda_F = \sum_{(\delta, a) \,\in\, {\rm terms}(F,m)} {\rm mean}(a) \; \ge \; 0,
\label{eq:svdx_lambda}
\end{equation}
where the mean is taken over all entries of the term array $a$. A low-overlap edge channel shows
up as a small value of $\lambda_F$. Since the entries of $a$ are non-negative, $\lambda_F = 0$
happens only if channel $F$'s contribution vanishes identically; the code drops such channels
rather than dividing by zero (their columns of $\bar A$, and rows of $U_{\rm freq}$ below, are
zero).

Now group the (channel, term) pairs by $\delta$. Since a {\tt PfVariance}'s terms are keyed by
$\delta$, a channel contributes at most one term to each group -- but a two-term channel appears
in two groups, which is why the groups' channel sets can overlap. For each $\delta$, stack the
flattened term arrays as the {\em columns} of a matrix $A_\delta$ of shape
$(2^{{\rm popc}(\delta)} P,\ n_{m\delta})$, dividing the column of channel $F$ by $\lambda_F$,
where $n_{m\delta}$ is the number of (channel, term) pairs in the group, and let $v_\delta,
\lambda_\delta$ denote the corresponding slices of $v, \lambda$. Then
\begin{equation}
y_m = \sum_\delta E_\delta \, A_\delta \, (v_\delta \odot \lambda_\delta),
\label{eq:svdx_group_sum}
\end{equation}
where $\odot$ denotes elementwise multiplication, and $E_\delta$ is the
$(DP) \times (2^{{\rm popc}(\delta)} P)$ zero/one ``expansion'' matrix which replicates the
packed DM axis over the unselected bits, i.e.\ copies packed entry $({\rm sel}(d,\delta), p)$ to
entry $(d, p)$. (In the code, $E_\delta$ is an index gather, not a materialized matrix.)

\subsection{Per-group SVD and stacking}
\label{ssec:svd_exact_stacking}

Each $A_\delta$ turns out to be numerically low-rank (variances vary smoothly with channel).
We compute the truncated SVD
\begin{equation}
A_\delta = W_\delta \, {\rm diag}(S_\delta) \, Q_\delta^T,
\end{equation}
where $W_\delta$ is $(2^{{\rm popc}(\delta)}P) \times K_\delta$ and $Q_\delta$ is
$n_{m\delta} \times K_\delta$, both with orthonormal columns, keeping the $K_\delta$ singular
values with $S > \varepsilon S_{\rm max}$ (\S\ref{ssec:svd_exact_threshold}).

Let $K_{\rm tot} = \sum_\delta K_\delta$. Substituting the truncated SVDs into
Eq.\ (\ref{eq:svdx_group_sum}) and stacking the per-group factors along the rank axis, we obtain
\begin{equation}
\bar A = \bar W \, {\rm diag}(\bar S) \, \bar Q^T,
\label{eq:svdx_stacked}
\end{equation}
where:
\begin{itemize}
\item $\bar W$ ($DP \times K_{\rm tot}$) concatenates the expanded left factors
$E_\delta W_\delta$ horizontally;
\item $\bar S$ (length $K_{\rm tot}$) concatenates the $S_\delta$;
\item $\bar Q$ ($n_m \times K_{\rm tot}$) concatenates the $Q_\delta$ horizontally, with each
group's rows scattered into its channel positions (zero-padded elsewhere).
\end{itemize}
This eliminates the sum over $\delta$, but the columns of $\bar W, \bar Q$ are {\em not}
orthonormal ($\bar Q^T \bar Q \ne {\rm id}$): different groups' channel sets overlap (via the
two-term channels), and the expanded columns $E_\delta W_\delta$ are not orthogonal across
groups -- nor even unit norm, since
$\|E_\delta w\|^2 = 2^{\rho-{\rm popc}(\delta)}\|w\|^2$. So $\bar S$ is not yet a set of
singular values of anything.

\subsection{Canonical form}
\label{ssec:svd_exact_canonical}

Finally, we reduce the stacked representation to a genuine truncated SVD of $\bar A$, without
materializing any $\bar A$-sized matrix. Take thin SVDs of the two stacked factors:
\begin{equation}
\bar W = W_1 \, {\rm diag}(W_2) \, W_3^T, \hspace{1.5cm} \bar Q = Q_1 \, {\rm diag}(Q_2) \, Q_3^T,
\end{equation}
where $W_1$ ($DP \times k_W$), $Q_1$ ($n_m \times k_Q$), $W_3$ ($K_{\rm tot} \times k_W$) and
$Q_3$ ($K_{\rm tot} \times k_Q$) have orthonormal columns, with
$k_W = \min(DP,\, K_{\rm tot})$ and $k_Q = \min(n_m,\, K_{\rm tot})$. Then
\begin{equation}
\bar A = W_1 \, M \, Q_1^T,
\hspace{1.5cm}
M \equiv {\rm diag}(W_2) \, W_3^T \, {\rm diag}(\bar S) \, Q_3 \, {\rm diag}(Q_2),
\label{eq:svdx_M}
\end{equation}
where $M$ is a small $(k_W \times k_Q)$ matrix. Eq.\ (\ref{eq:svdx_M}) is an identity, not an
approximation -- the only error so far is the per-group truncation. It holds even if $\bar W$ or
$\bar Q$ is rank deficient: a near-null direction shows up as a roundoff-sized entry of $W_2$ or
$Q_2$, which multiplies the corresponding row or column of $M$ and thereby kills the (arbitrary)
direction in $W_1$ or $Q_1$. Using $\min$ shapes means that no assumption such as
$K_{\rm tot} \le n_m$ is needed.

Since the columns of $W_1, Q_1$ are orthonormal, the singular values of $M$ are exactly the
singular values of $\bar A$. SVD-decompose $M = X \, {\rm diag}(S_M) \, Y^T$, let $K_0$ be the
number of singular values above the same relative threshold $\varepsilon$, and absorb the
factors into $W_1, Q_1$:
\begin{equation}
S_0 = (S_M)_{[:K_0]}, \hspace{1cm}
U_{\rm out} = W_1 \, X_{[:,\,:K_0]}, \hspace{1cm}
U_{\rm freq} = Q_1 \, Y_{[:,\,:K_0]},
\end{equation}
where $[:,:K_0]$ denotes the first $K_0$ columns. This gives the canonical representation
\begin{equation}
\bar A = U_{\rm out} \, {\rm diag}(S_0) \, U_{\rm freq}^T + O(\varepsilon \|\bar A\|_2),
\hspace{1.5cm}
U_{\rm out}^T U_{\rm out} = U_{\rm freq}^T U_{\rm freq} = {\rm id}_{K_0},
\label{eq:svdx_canonical}
\end{equation}
with $S_0$ sorted descending -- i.e.\ a genuine truncated SVD of $\bar A$, with $K_0$ its
numerical rank. The three factors are the {\tt VarianceMapExact} block members {\tt sv\_out},
{\tt sval} and {\tt sv\_freq}, stored with the rank axis first.

Given an updated $v$, reevaluating the variances is two small matrix-vector multiplies and two
diagonal scalings,
\begin{equation}
y_m = U_{\rm out} \, \big( S_0 \odot ( U_{\rm freq}^T \, \Lambda \, v ) \big),
\label{eq:svdx_eval}
\end{equation}
i.e.\ $O\big(K_0 (n_m + DP)\big)$ flops per $(\mbox{tree}, m)$; likewise the storage (the
factors $U_{\rm out}, S_0, U_{\rm freq}$ plus the length-$n_m$ vector $\lambda$) is
$O\big(K_0 (n_m + DP)\big)$, instead of $O(n_m \cdot DP)$ for a dense $\bar A$.

\subsection{Truncation threshold}
\label{ssec:svd_exact_threshold}

Truncating at $\varepsilon S_{\rm max}$ perturbs the reconstruction at the $\varepsilon$ level
relative to $\|\bar A\|_2$, so $\varepsilon$ is an accuracy knob -- but it is only meaningful
above the double-precision noise floor on singular values, which numpy's {\tt matrix\_rank}
estimates as $\max(\mbox{shape}) \cdot \epsilon_{\rm f64} \cdot S_{\rm max}$.
A fixed threshold is therefore not safe in general. An earlier prototype of this construction
used $\varepsilon = 10^{-11}$, which sits above the floor at the sizes it was run at, but the exact
map's output axis $DP = 2^{r-R}P$ is larger than the approximation's $2^{r-L}P$ by $2^{L-R}$,
and at $DP = 65536$ the floor is $\approx 1.5 \times 10^{-11}$, i.e.\ {\em above} $10^{-11}$. We
therefore choose $\varepsilon$ per block from its largest matrix dimension,
\begin{equation}
\varepsilon = \max\Big( 10^{-11}, \;\; 16 \cdot \max(DP,\, n_m) \cdot \epsilon_{\rm f64} \Big),
\label{eq:svdx_epsilon}
\end{equation}
and apply it to every SVD in that block. The factor 16 sits comfortably above the estimated
floor while remaining far below any genuine component. (Both {\tt VarianceMapExact} and
{\tt VarianceMapApproximation} take an {\tt epsilon} constructor argument which overrides this
default.)

In practice the choice turns out not to matter much, because the spectra have clean cliffs:
on {\tt chord\_sb2\_et.yml}, where the two rules differ by a factor $\sim\!6$
($5.8\times10^{-11}$ against $10^{-11}$), the approximate map's total rank
$\sum K_0 = 4969$ is {\em identical} under both, and $\sum K_{\rm tot}$ differs by 2 out of
$\sim\!11300$. The shape-aware rule is used because it degrades gracefully as the matrices grow,
not because it changes today's answers.

That $K_0$ is a genuine numerical rank, rather than an artifact of where the threshold is put,
can be checked directly: for a representative {\tt toy.yml} block ($n_m = 640$,
$DP = 6656$), $K_0 = 175$ or $176$ for every $\varepsilon$ from $10^{-8}$ down to $10^{-12}$ --
the spectrum has a clean cliff, with the last retained singular value at
$4 \times 10^{-11} S_{\rm max}$ and the next one below $10^{-12} S_{\rm max}$. Loosening to
$\varepsilon = 10^{-6}$ gives $K_0 = 135$.

\subsection{Cost, storage and validation}
\label{ssec:svd_exact_cost}

For {\tt toy.yml} ($N_{\rm freq} = 640$, 6 trees, 110 $(\mbox{tree},m)$ pairs), we find
$n_m$ from $159$ to $640$, $DP$ from $896$ to $6656$, $K_{\rm tot}$ from $90$ to $294$ (summing
to $18131$), and $K_0$ from $45$ to $178$ (summing to $9812$) -- so the canonical form buys a
further factor $\sim\!1.9$ beyond the per-group SVDs. Total storage is $261$ MiB, versus $907$
MiB for the dense $\bar A$'s, a factor $3.5$. Building {\tt PfAvarExact} takes $3.5$ sec and
building {\tt VarianceMapExact} from it takes $11$ sec.

The compression factor is modest here only because {\tt toy.yml} is small in frequency: when
$DP \gg n_m$, Eq.\ (\ref{eq:svdx_eval}) gives a storage ratio of $n_m DP / K_0(n_m + DP) \to
n_m / K_0$, which is $\sim\!3.6$ at $n_m = 640$. Configs with more input channels compress
better, since $K_0$ grows much more slowly than $n_m$.

Reconstructing the variances through Eq.\ (\ref{eq:svdx_eval}) agrees with the directly-computed
{\tt PfAvarExact} values to a relative accuracy of $4 \times 10^{-12}$ on {\tt toy.yml}, and
$\sim\!10^{-14}$ on small random configs, consistent with the threshold $\varepsilon$. The check
that matters is done at an input vector $v$ which {\tt PfAvarExact} was {\em not} built with,
recomputing the reference directly from {\tt per\_tfm}: this tests that we have captured the
linear map, not merely its value at one input. It is available as
{\tt VarianceMapExact.check()}.

Finally, a note on scale. At {\tt chord\_sb2\_et.yml} parameters ($r = 16$, $R = 4$,
$DP = 65536$, $M = 91$, 10 trees, so $\sim\!910$ blocks), $U_{\rm out}$ alone would run to tens
of GB. This is moot, however, since {\tt PfAvarExact} cannot be built at that scale in the first
place (\S\ref{sec:avar_pfvariance}). Should storage become the binding constraint at some
intermediate scale, the lever is to stop materializing the expanded output-side factor: writing
$U_{\rm out} = \sum_\delta E_\delta W_\delta C_\delta$ for small $C_\delta$ of shape
$K_\delta \times K_0$, a block would store the packed $W_\delta$ (with $2^{{\rm popc}(\delta)}P$
rows) and the $C_\delta$, instead of $K_0 \cdot DP$ doubles.

\clearpage

\section{{\tt PfAvarApproximation}, {\tt VarianceMapApproximation}}
\label{sec:variance_svd}

This appendix is the approximate counterpart of \S\ref{sec:analytic_variance}, and assumes it
has been read: {\tt PfAvarApproximation} computes the peak-finding variances $y = Av$ for one
input vector $v$, and {\tt VarianceMapApproximation} compresses the map $A$ itself into a
canonical low-rank form. Both are in \verb|pirate_frb/slow_avar/|. The linear algebra is
identical to \S\ref{ssec:svd_exact_stacking}--\S\ref{ssec:svd_exact_canonical} and is not
repeated here; what follows is the delta.

The reason both pairs of classes exist is cost. {\tt PfAvarExact} is only practical at moderate
rank (\S\ref{sec:avar_pfvariance}), so production uses the approximation described below, which
is also ported to C++ (\verb|pirate_frb/fast_avar/|). Correspondingly,
{\tt VarianceMapApproximation} {\em can} be run at CHIME/CHORD scale, and
{\tt VarianceMapExact} cannot.

\subsection{The approximation}
\label{ssec:approx_def}

Let $2^L$ be the tree's DM downsampling factor for peak-finding weights, where
$0 \le R \le L \le r$.\footnote{Code names: $2^L = {\tt wt\_dm\_downsampling}$,
$L = {\tt tree\_L}$.} Divide the tree-freqs into the $2^R$ {\em level-0 bands} (the aligned
blocks of $2^{r-R}$ tree-freqs appearing as subbands $(0,s)$ in \ddnotes{Subband search}), and
{\bf neglect the frequency-channel overlaps between them}. Two consequences:
\begin{itemize}
\item The variance depends on the DM index only through its $(r-L)$ highest bits -- the same DM
resolution as the peak-finding weights array. Writing a multiplet as (subband $n$, fine DM), the
variance then depends on $n$ but not on the fine DM.
\item The per-tree variance can be computed once as a $(2^{r-L}, 2^R, P)$ array and reduced
afterwards, either to $(2^{r-L}, N, P)$ by averaging the level-0 bands covered by each subband,
or to $(2^{r-L}, M, P)$ by the one-to-many subband-to-multiplet map.
\end{itemize}
So the natural object is the per-band variance array $y_j[d,p]$, $0 \le j < 2^R$,
$0 \le d < 2^{r-L}$, and {\tt tree\_variance} is the per-subband average
\begin{equation}
Y_n \;=\; \frac{1}{f_{\rm hi}(n) - f_{\rm lo}(n)} \sum_{j\,=\,f_{\rm lo}(n)}^{f_{\rm hi}(n)-1} y_j ,
\label{eq:approx_subband_avg}
\end{equation}
where $[f_{\rm lo}(n), f_{\rm hi}(n))$ is subband $n$'s coarse-freq range
({\tt n\_to\_flo}, {\tt n\_to\_fhi}).

The size of the error this introduces is not estimated here; it is measured directly, by
comparing {\tt PfAvarExact} and {\tt PfAvarApproximation} tree by tree
(\verb|pirate_frb check_avar_approximation|, which reports the distribution of
$\epsilon = {\rm var}_{\rm approx}/{\rm var}_{\rm exact} - 1$).

\subsection{Computing it: reuse across trees}
\label{ssec:approx_algorithm}

The approximation also makes the calculation much cheaper than
\S\ref{sec:avar_sparsetile}, in two ways.

First, the expensive part of the exact algorithm is gone. There, the cost was dominated by the
few input channels whose gridding footprint straddles a high-level coarse-freq boundary, whose
final tile has $O(2^{i_0})$ stored rows of length $O(2^{i_0})$. Neglecting inter-band overlaps
means we never merge across those boundaries: we stop iterating at the level where a coarse-freq
channel {\em is} a level-0 band, and read off one {\tt PfVariance} per (channel, band).

Second, the work is shared across trees. For each input channel we iterate a single
{\tt SparseTileTriple} through $k = 0, 1, 2, \ldots$, and a tree consumes it at its own
{\em k-level} $k = r - L + [\,{\rm ipri} > 0\,]$, at which point a singleton tile is converted to
a {\tt PfVariance} (rank $r-L$) and accumulated. Since many trees share a k-level, the loop is
ordered with $k$ outermost and the trees at that $k$ innermost, so the tile iteration and the
tile-to-variance conversion are each done once and reused. Trees with different $P$ share the
conversion by building it at the largest $P$ at that k-level and truncating on accumulation; a
downsampled tree (${\rm ipri} > 0$) takes the upper DM-half of the shared result.
The accumulation itself is the same two-array structure as \S\ref{sec:avar_pfvariance}:
{\tt per\_tff[t][ifreq][j]} holds one channel's contribution to band $j$, and
{\tt per\_tf[t][j]} its $v$-weighted sum over channels.

\subsection{The variance map, per level-0 band}
\label{ssec:approx_map}

The compression of \S\ref{ssec:svd_exact_map}--\S\ref{ssec:svd_exact_cost} now applies almost
verbatim, under two substitutions:
\begin{center}
\begin{tabular}{ll}
\S\ref{sec:analytic_variance} (exact) & this appendix (approximate) \\
\hline
block index: multiplet $m$            & block index: level-0 band $j$, $0 \le j < 2^R$ \\
DM bits $\rho = r-R$                  & DM bits $\rho = r-L$ \\
\end{tabular}
\end{center}
So a block is the $(2^{r-L}P) \times n_j$ matrix $A^{(j)}$ with $y_j = A^{(j)} v$, where the
contributing channels again form a contiguous range $F_{\rm lo} \le F < F_{\rm lo} + n_j$
(same argument as \S\ref{ssec:svd_exact_map}; asserted in the code).\footnote{Code names:
$F_{\rm lo} = {\tt per\_tf\_ifreq\_lo}$, $n_j = {\tt per\_tf\_nf}$, members of
{\tt PfAvarApproximation}.} Everything from there on -- the preconditioner
$\Lambda$, the grouping by {\tt dbits}, the per-group SVDs, the stacking, the canonical form
$\bar A = U_{\rm out}\,{\rm diag}(S_0)\,U_{\rm freq}^T$, and the threshold
(\ref{eq:svdx_epsilon}) -- is unchanged.

The one new ingredient is Eq.\ (\ref{eq:approx_subband_avg}). It is linear, so it composes
trivially with the blocks: {\tt eval\_tree()} evaluates each of the $2^R$ blocks once and then
averages. This is also why the block index is the band $j$ and not the subband $n$: a band
belongs to several subbands, and there can be more subbands than bands (for
{\tt chord\_sb2\_et.yml} tree 0, $N = 25$ against $2^R = 16$), so blocking by $n$ would repeat
work.

\subsection{Results}
\label{ssec:approx_results}

For {\tt chord\_sb2\_et.yml} (10 trees, 136 $(\mbox{tree}, j)$ blocks), the python
{\tt PfAvarApproximation} builds in 13 sec and {\tt VarianceMapApproximation} in a further
13 sec. Almost all channels produce exactly one term: 213 of the 174496 (channel, band) pairs
produce two, the rest one. The number of distinct bitmasks $\delta$ per block runs from 7 to 32.
We find $n_j$ from $270$ to $2823$, $2^{r-L}P$ from $448$ to $16384$, $K_{\rm tot}$ from $20$ to
$160$ (summing to $11299$), and $K_0$ from $12$ to $69$ (summing to $4969$) -- so the canonical
form buys a further factor $\sim\!2.3$ beyond the per-group SVDs. Total storage is $241$ MiB
against $7807$ MiB for the dense $\bar A$'s, a factor $32$.

That compression factor is much better than the $3.5$ of \S\ref{ssec:svd_exact_cost}, and for
the reason given there: the ratio tends to $n_j/K_0$ when $2^{r-L}P \gg n_j$, and $n_j$ here
runs into the thousands rather than being capped at {\tt toy.yml}'s 640. Reconstructing the
variances through the canonical form agrees with the directly-computed {\tt PfAvarApproximation}
values to a relative accuracy of $6 \times 10^{-11}$, consistent with $\varepsilon$; as in
\S\ref{ssec:svd_exact_cost}, the check that matters is run at an input vector which the
{\tt PfAvarApproximation} was not built with ({\tt VarianceMapApproximation.check()}).

Two practical notes. The C++ {\tt PfAvarApproximation} keeps only {\tt per\_tf} and
{\tt tree\_variance}, dropping the per-channel {\tt per\_tff} which the compression needs, so
{\tt VarianceMapApproximation} must be built from the python class; at 13 sec this is not
currently a problem. And none of this is yet wired into the real-time pipeline.

\end{document}
